\documentclass[11pt,oneside]{book}
\usepackage[margin=1.2in]{geometry}
\usepackage[toc,page]{appendix}
\usepackage{graphicx}
\usepackage{natbib}
\usepackage{lipsum}
\usepackage{caption}
\usepackage[T1]{fontenc}
\usepackage{titlesec, blindtext, color}
\usepackage{xcolor,tikz}
\usetikzlibrary{patterns}
\usetikzlibrary{decorations.markings}
\usepackage{amsmath,amssymb,amsthm,mathrsfs,amsfonts,xfrac,pifont,bbold,physics}
\usepackage[utf8]{inputenc}
\usepackage{amsthm}
\usepackage[breakable, theorems, skins]{tcolorbox}
\usepackage[colorlinks = true,
            linkcolor = red,
            urlcolor  = blue,
            citecolor = red,
            anchorcolor = red]{hyperref}
\usepackage{cleveref}

\usepackage{soul}
\usepackage{frcursive}
\usepackage{booktabs}

\usepackage{array}
\newcolumntype{C}[1]{>{\centering\arraybackslash}m{#1}}

% -------------------------------------------------------------------
% Theorem Styles
% -------------------------------------------------------------------

\theoremstyle{definition} % Define theorem styles here based on the definition style (used for definitions and examples)
\newtheorem*{definition}{Definition}

\theoremstyle{plain} % Define theorem styles here based on the plain style (used for theorems, lemmas, propositions)
\newtheorem{theorem}{Theorem}[section]
\newtheorem{axiom}{Axiom}
\newtheorem{corollary}[theorem]{Corollary}
\newtheorem{lemma}[theorem]{Lemma}
\newtheorem{proposition}[theorem]{Proposition}

\theoremstyle{remark} % Define theorem styles here based on the remark style (used for remarks and notes)
\newtheorem*{notation}{Notation}
\newtheorem*{solution}{Solution}

\newtheoremstyle{underline}% name
{}        % Space above, empty = `usual value'
{}              % Space below
{}              % Body font
{}    % Indent amount (empty = no indent, \parindent = para indent)
{}              % Thm head font
{.}             % Punctuation after thm head
{1.5mm}         % Space after thm head: \newline = linebreak
{{\underline{\textit{\thmname{#1}\thmnumber{ #2}}~\thmnote{(#3)}\unskip}}}  % Thm head spec

\theoremstyle{underline}

\newtheorem{remark}[theorem]{Remark}
\newtheorem{example}[theorem]{Example}
\newtheorem{claim}[theorem]{Claim}
% -------------------------------------------------------------------
% Chapter Headings
% -------------------------------------------------------------------

%\setcounter{chapter}{-1}

\makeatletter
\renewcommand{\@chapapp}{Lecture}
\makeatother
\definecolor{lightergray}{rgb}{0.9,0.9,0.9}

\usepackage{titlesec}
\titleformat{\section}{\large\bfseries\raggedright}{}{0em}{\colorsection}[\titlerule]
\titleformat{name=\section,numberless}{\large\scshape\bfseries\raggedright}{}{0em}{\colorsectionnonumber}[\titlerule]

\titleformat{\subsection}{\bfseries\raggedright}{}{0em}{\colorsubsection}
\titleformat{name=\subsection,numberless}{\bfseries\raggedright}{}{0em}{\colorsubsectionnonumber}

\newcommand{\colorsection}[1]{%
    \colorbox{lightergray}{\parbox{\dimexpr\textwidth-2\fboxsep}{\thesection\ \ #1}}}
\newcommand{\colorsectionnonumber}[1]{%
    \colorbox{lightergray}{\parbox{\dimexpr\textwidth-2\fboxsep}{#1}}}
    
\newcommand{\colorsubsection}[1]{%
    \colorbox{lightergray}{\parbox{\dimexpr\textwidth-2\fboxsep}{\thesubsection\ #1}}}
\newcommand{\colorsubsectionnonumber}[1]{%
    \colorbox{lightergray}{\parbox{\dimexpr\textwidth-2\fboxsep}{#1}}}
    
\definecolor{gray75}{gray}{0.75}
\newcommand{\hsp}{\hspace{20pt}}
\titleformat{\chapter}[hang]{\Huge\bfseries}{\thechapter\hsp\textcolor{gray75}{|}\hsp}{0pt}{\Huge\bfseries}

% Begin/end Shortcuts
\def\bc{\begin{corollary}}
\def\ec{\end{corollary}}
\def\bcl{\begin{claim}}
\def\ecl{\end{claim}}
\def\bd{\begin{definition}}
\def\ed{\end{definition}}
\def\ben{\begin{enumerate}}
\def\een{\end{enumerate}}
\def\be{\begin{equation}}
\def\ee{\end{equation}}
\def\bse{\begin{equation*}}
\def\ese{\end{equation*}}
\def\bex{\begin{example}}
\def\eex{\end{example}}
\def\bit{\begin{itemize}}
\def\eit{\end{itemize}}
\def\bl{\begin{lemma}}
\def\el{\end{lemma}}
\def\bnn{\begin{notation}}
\def\enn{\end{notation}}
\def\bn{\begin{note}}
\def\en{\end{note}}
\def\bp{\begin{proposition}}
\def\ep{\end{proposition}}
\def\bq{\begin{proof}}
\def\eq{\end{proof}}
\def\br{\begin{remark}}
\def\er{\end{remark}}
\def\bs{\begin{solution}}
\def\es{\end{solution}}
\def\btab{\begin{table}}
\def\etab{\end{table}}
\def\btb{\begin{tabular}}
\def\etb{\end{tabular}}
\def\bt{\begin{theorem}}
\def\et{\end{theorem}}
\def\btik{\begin{tikzpicture}}
\def\etik{\end{tikzpicture}}

% Letter Shortcuts
% Greek letters
\def\a{\alpha}
\def\b{\beta}
\def\del{\delta}
\def\g{\gamma}
\def\l{\lambda}
\def\n{\nabla}
\def\sig{\sigma}

% Mathbb
\def\b1{\mathbb{1}} % identity
\def\C{\mathbb{C}} % complex numbers
\def\F{\mathbb{F}}
\def\N{\mathbb{N}} % Natural
\def\Q{\mathbb{Q}} % Rational
\def\R{\mathbb{R}} % Reals
\def\Z{\mathbb{Z}} % Integers

% Mathcal
\def\cA{\mathcal{A}}
\def\cB{\mathcal{B}}
\def\cD{\mathcal{D}}
\def\cF{\mathcal{F}}
\def\cH{\mathcal{H}}
\def\cI{\mathcal{I}}
\def\cJ{\mathcal{J}}
\def\cK{\mathcal{K}}
\def\cL{\mathcal{L}}
\def\cM{\mathcal{M}}
\def\cN{\mathcal{N}}
\def\cO{\mathcal{O}}
\def\cP{\mathcal{P}}
\def\cT{\mathcal{T}}

% Mathfrak 
\DeclareMathOperator{\GL}{GL}
\def\gl{\mathfrak{gl}}
\DeclareMathOperator{\Ort}{O}
\def\ort{\mathfrak{o}}
\DeclareMathOperator{\SL}{SL}
\def\sl{\mathfrak{sl}}
\DeclareMathOperator{\SO}{SO}
\def\so{\mathfrak{so}}
\DeclareMathOperator{\SU}{SU}
\def\su{\mathfrak{su}}
\def\fF{\mathfrak{F}}

% Misc. 
\def\diag{\text{diag}}
\def\p{\partial}
\def\ss{\subset}
\def\se{\subseteq}
\def\sm{\setminus}
\def\ve{\varepsilon}
\def\vn{\varnothing}
\def\cl{\colon}

% Boxes 
\DeclareRobustCommand{\mybox}[2][colframe=blue!10!black,before skip=10pt]{%
\begin{tcolorbox}[
        breakable,
        width=\dimexpr\textwidth\relax, 
        boxsep=5pt,
        arc=5pt,outer arc=5pt,
        ]
        #2
\end{tcolorbox}
}

\def\bbox{\begin{tcolorbox}[colback=white!5,colframe=blue!75!black,title=Exercise]}
\def\ebox{\end{tcolorbox}}

% Tikz Stuff
\def\midarrow{\draw[thick, decoration={markings, mark=at position 0.5 with {\arrow{>}}}, postaction={decorate}]}
\def\wavey{\draw[thick, decorate, decoration={snake, segment length=1.5mm, amplitude=0.5mm}]}

\begin{document}

\captionsetup[figure]{margin=1.5cm,font=small,labelfont={bf},name={Figure},labelsep=colon,textfont={it}}
\captionsetup[table]{margin=1.5cm,font=small,labelfont={bf},name={Table},labelsep=colon,textfont={it}}
\setlipsumdefault{1}

\frontmatter

\begin{titlepage}
	\centering
	    \scshape % Use small caps for all text on the title page
        \vspace*{\baselineskip} % White space at the top of the page
        
	    \rule{\textwidth}{1.6pt}\vspace*{-\baselineskip}\vspace*{2pt} % Thick horizontal rule
	    \rule{\textwidth}{0.4pt} % Thin horizontal rule
	    
	    \vspace{0.75\baselineskip} % Whitespace above the title
	    
	    {\LARGE QFT I (Path Integrals)} % Title
	    
	    \vspace{0.75\baselineskip} % Whitespace below the title
	    
	    \rule{\textwidth}{0.4pt}\vspace*{-\baselineskip}\vspace{3.2pt} % Thin horizontal rule
	    \rule{\textwidth}{1.6pt} % Thick horizontal rule
	    
        \vspace{5\baselineskip} % Whitespace after the title block

	    Course delivered in 2019 by 
	
	    \vspace{0.5\baselineskip} % Whitespace before the editors
	
	    {\scshape\Large Dr. Douglas Smith} % Lecturer Name
	
	    \vspace{0.5\baselineskip} % Whitespace below the editor list
	
	    \textit{Durham University} % Lecturer Institution
	    
	    \vspace{2\baselineskip} % Whitespace after the title block
	    
	    \includegraphics[width=8cm]{images/DurhamLogo.png}\\[1cm] % Logo 

	    Notes taken by 
	
	    \vspace{0.5\baselineskip} % Whitespace before my name
	
	    {\scshape\Large Richie Dadhley} % Lecturer Name
	   
	    \vspace{0.5\baselineskip} % Whitespace below my name
	    \textit{richie@dadhley.com} % Email
	
	    \vfill % Whitespace between editor names and publisher logo
\end{titlepage}



% -------------------------------------------------------------------
% Acknowledgements
% -------------------------------------------------------------------

\newpage
\section*{Acknowledgements}

These are my notes on the 2019 lecture course "QFT I" taught by Dr. Douglas Smith at Durham University as part of the Particles, Strings and Cosmology Msc. For reference, the course lasted 8 hours and was taught over the course of 4 weeks. This course is really meant more as an introduction to path integrals then a study of QFT. The module QFT II will do the actual QFT. \\

I have tried to correct any typos and/or mistakes I think I have noticed over the course. I have also tried to include additional information that I think supports the taught material well, which sometimes has resulted in modifying the order the material was taught. Obviously, any mistakes made because of either of these points are entirely mine and should not reflect on the taught material in any way. \\

% For any other notable sources used:
% One of the main sources I have used to obtain the additional information is Dr. David Tong’s String Theory notes, available via the following link  

%\begin{center}
%    \textcolor{red}{Link}
%\end{center}

I would like to extend a message of thanks to Dr. Smith for teaching this course. \\

If you have any comments and/or questions please feel free to contact me via the email provided on the title page. \\

For a list of other notes/works I have available, visit my blog site

\begin{center}
    \href{https://richie291.wixsite.com/theoreticalphysics}{https://richie291.wixsite.com/theoreticalphysics}
\end{center}

These notes are not endorsed by Dr. Smith or Durham University.

\vspace{1cm}

\begin{flushright}
    \Huge{{\cursive\setul{0.1ex}{}\ul{~~Richie Dadhley~~}}}
\end{flushright}

% -------------------------------------------------------------------
% Contents
% -------------------------------------------------------------------

\tableofcontents

% -------------------------------------------------------------------
% Main sections (as required)
% -------------------------------------------------------------------

\mainmatter

%\chapter{Introduction}


Lecture one - Wick rotation end  (skip Greens functions)

lecture 2 - up to Feynman kernel (classical limit next lecture) 

lecture 3 - derived SD equation 

lecture 4 - rest 
\chapter{Path Integrals \& Feynman Kernel}

This course is the study of path integrals in quantum mechanics (QM), with a particular emphasis on concepts related to quantum field theory (QFT). The first obvious questions to ask are "what are path integrals?" and "why do we care?" The first question will be answered as we go along, but allow us to answer the second now, providing motivation for the course. 

\section{Why Do We Need Path Integrals?}

When we are first taught QM it is almost always taught\footnote{If you were taught Path integral QM first, fair play.} in the language of Schr\"{o}dinger wavefunctions and Dirac bra-ket notation. The central equation is the Schr\"{o}dinger equation
\bse 
    i\frac{d}{dt}\ket{\psi} = \hat{H}\ket{\psi},
\ese 
where $\hat{H}$ is the operator version of the classical Hamiltonian. As such it is clear that this approach to QM is based on the \textit{Hamiltonian} formulation of classical mechanics.

A person who is familiar with classical field theory, will know that there is a completely different way to study classical mechanics; in terms of the \textit{Lagrangian}. The path integral formulation approaches QM from this perspective. That is the Schr\"{o}dinger equation does not appear, but instead we study principles of least action and partition functions. 

Now the latter generally involves \textit{a lot} of integral manipulation,\footnote{Trust me, you'll agree shortly.} and so you would be justified in asking "why on Earth would we put ourselves through all this pain when we could just use the Hamiltonian approach to QM?" The answer to that question is the name of this course... QFT. If you are reading these notes I think it's fair to assume you have had some introduction to QFT, and it was most definitely done in the language of so-called \textit{second quantisation}. This is the Hamiltonian approach to QFT and is obviously a very powerful tool. However, you will probably have been told about the perils of breaking manifest Lorentz invariance\footnote{For clarity, by "manifest" Lorentz invariant we mean you can just look at the equations and say "yup that's Lorentz invariant."} in this approach, and then having to make checks along the way to fix your abuse. For example, you have to introduce a factor of $1/\sqrt{2E_{\Vec{p}}}$ in order to get a Lorentz invariant measure, which then means you need to renormalise your states. 

Of course this all works, but it would be nice if we could formulate the study of QFT in a manifestly Lorentz invariant way. This is where the path integral formulation comes in. Lagrangians are essentially \textit{defined} to be manifestly Lorentz invariant, so if we make them our central objects, we're good to go. I believe this is motivation enough for you to believe that its at least worth studying. So let's get to it.

\br 
    It is worth clarifying, it is not that the Hamiltonian doesn't appear in the path integral approach, its that we do not make it the central object. It is the Lagrangian that we are fundamentally interested in. 
\er 

\section{What Do We Want To Calculate?}

For simplicity, in this course we will just focus on 1-dimensional mechanics, i.e. a particle with position $q(t)$ and momentum $p(t)$, with standard kinetic energy 
\bse 
    \frac{m\dot{q}^2}{2} = \frac{p^2}{2m},
\ese 
and potential $V(q)$. The generalisation to higher dimensions is not hard to see, but involves carrying around a bunch of sums/dot products, and as we shall see there's enough symbols as is so let's not make things worse. We therefore have a time-independent Hamiltonian 
\be
\label{eqn:Hamiltonian}
    \hat{H} = \frac{\hat{q}^2}{2m} + \hat{V}(\hat{q})
\ee 
and action 
\be
\label{eqn:ActionAndLagrangian}
    S = \int dt \, L(q,\dot{q}), \qquad L(q,\dot{q}) = \frac{m}{2}\dot{q}^2 - V(q) . 
\ee 

The basic object we will calculate using path integrals is the \textit{Feynman kernel}
\mybox{
    \be 
    \label{eqn:FeynmanKernel}
        K(q_F,t_F;q_I,t_I) := \bra{q_F} \hat{U}(t_F,t_I) \ket{q_I}, 
    \ee 
}
\noindent where
\be 
\label{eqn:TimeEvolutionOperator}
    \hat{U}(t,t_0) := \exp\bigg(-\frac{i}{\hbar} \hat{H}(t-t_0)\bigg)
\ee 
is the \textit{time evolution operator}. The Feynman kernel tells us the amplitude for a particle that is initially at position $q_I$ at time $t_I$ to evolve to position $q_F$ at time $t_F$, where of course $t_F>t_I$. As such, the Feynman kernel contains all the information about the dynamics of the system, which is why we care about it so much.

\br 
    Note that for a time-independent Hamiltonian, we have 
    \bse 
        \hat{U}(t,t_0) = \hat{U}(T) = \exp\bigg(-\frac{i}{\hbar}\hat{H}T\bigg),
    \ese
    where $T=t-t_0$. So the Feynman kernel becomes 
    \bse 
        K(q_f,t_F;q_I,t_I) = K(q_F,T;q_I,0),
    \ese 
    where $T=t_F-t_I$. We will use this from now on, unless otherwise specified.
\er 

The trick to getting the Feynman integral is to insert a complete set of states
\bse 
    \b1 = \int dq \bra{q}\ket{q}
\ese 
for both the initial and final state. Let's give examples for clarity.

\bex 
    The amplitude for some initial state $\ket{\psi}$ to evolve to some final state $\ket{\phi}$ in time $T$ is 
    \bse 
        \begin{split}
            \bra{\phi} \hat{U}(T) \ket{\psi} & = \int dq_I dq_F \braket{\phi}{q_F} \bra{q_F} \hat{U}(T) \ket{q_I} \braket{q_I}{\psi} \\
            & = \int dq_I dq_F \braket{\phi}{q_F} K(q_F,T;q_I,0) \braket{q_I}{\psi} \\
            & = \int q_I q_F \, \phi^*(q_F) K(q_F,T;q_I,0) \psi(q_I),
        \end{split}
    \ese 
    where at the last line we've used the usual QM notation for the position space representation of a state:
    \bse 
        \psi(q) := \braket{q}{\psi}.
    \ese 
\eex

\bex 
    Let's assume we have some state with initial wavefunction $\psi(q_I,t_I)$. We can find the expression for its evolution to some later wavefunction $\psi(q_F,t_F)$ in terms of the Feynman kernel:
    \bse 
        \begin{split}
            \psi(q_F,t_F) & := \braket{q_F}{\psi(t_F} \\
            & = \bra{q_F} \hat{U}(T) \ket{\psi(t_I)} \\
            & = \int dq_I \bra{q_F} \hat{U}(T) \ket{q_I} \braket{q_I}{\psi(t_I)} \\
            & =\int dq_I K(q_F,T;q_I,0) \psi(q_I,t_I),
        \end{split}
    \ese 
    where we have used the position space representation definition given at the end of the last example and also the definition of the time evolution operator, 
    \bse 
        \ket{\psi(t_2)} = \hat{U}(t_2,t_1)\ket{\psi(t_1)},
    \ese 
    along with $T=t_F-t_I$.
\eex 

To give some foresight, we will also use the Feynman kernel to calculate what are called \textit{time-ordered correlations functions} or \textit{Green's functions}, 
\bse 
    G(t_1,t_2,...,t_N) := \bra{\Omega} \cT \big\{\hat{q}(t_1)\hat{q}(t_2) ... \hat{q}(t_N) \big\}\ket{\Omega},
\ese 
where $\ket{\Omega}$ is the ground state of our system, and $\cT$ is the time ordering operator, which essentially orders the Heisenberg\footnote{As otherwise the operators are time independent} picture operators in increasing time, e.g. 
\bse 
    \cT\big\{\hat{\cO}_1(t_1)\hat{\cO}_2(t_2)\big\} = \begin{cases}
        \hat{\cO}_1(t_1)\hat{\cO}_2(t_2) & \text{if } t_1>t_2 \\
        \hat{\cO}_2(t_2)\hat{\cO}_1(t_1) & \text{if } t_2>t_1
    \end{cases}.
\ese 

\bnn 
    From now on I am going to drop the hats on operators in order to save typing. It should be reasonably clear from the context, and whenever potential confusion might arise I shall try be explicit. 
\enn 

\section{Motivating The Path Integral Approach}

So what is a path integral? Well as the name suggests, it is an integral over paths, but what do we mean by this and what's it got to do with the Feynman kernel? Well let's consider the double slit experiment. The interference pattern this generated can be explained in terms of waves, or we can explain it in terms of particles, provided we put a phase on each particle path. That is, the total amplitude to go from the source to a particular point is given by the sum of the amplitudes of the two separate paths, as indicated in the figure below. 

\begin{center}
    \btik 
        \node at (0,0) {\Huge{$*$}};
        \node at (0,-0.5) {Source};
        \draw[thick] (2,2) -- (2,1.1);
        \draw[thick] (2,0.9) -- (2,-0.9);
        \draw[thick] (2,-1.1) -- (2,-2);
        \midarrow (0,0) -- (2,1);
        \midarrow (2,1) -- (6,0.5);
        \node at (4,1.2) {Path 1, $\cA_1$};
        \midarrow (0,0) -- (2,-1);
        \midarrow (2,-1) -- (6,0.5);
        \node at (4,-0.8) {Path 2, $\cA_2$};
        \draw[fill=black] (6,0.5) circle [radius=0.07] node [right] {$\cA=\cA_1+\cA_2$};
    \etik 
\end{center}

\subsection{Zee's "Wise Guy"}

The story then goes as follows:\footnote{I'm paraphrasing Zee's book here, as I love this description.} a teacher had just finished explaining the above idea to his class when a student, going by the name of Feynman, chimed in and said "Excuse me sir, but what happens if I put another slit in my screen?" The lecturer then replied "We just add another contribution" while drawing the following diagram

\begin{center}
    \btik 
        \node at (0,0) {\Huge{$*$}};
        \node at (0,-0.5) {Source};
        \draw[thick] (2,2) -- (2,1.1);
        \draw[thick] (2,0.9) -- (2,-0.4);
        \draw[thick] (2,-0.6) -- (2,-0.9);
        \draw[thick] (2,-1.1) -- (2,-2);
        \midarrow (0,0) -- (2,1);
        \midarrow (2,1) -- (6,0.5);
        \midarrow (0,0) -- (2,-1);
        \midarrow (2,-1) -- (6,0.5);
        \midarrow (0,0) -- (2,-0.5);
        \midarrow (2,-0.5) -- (6,0.5);
        \draw[fill=black] (6,0.5) circle [radius=0.07] node [right] {$\cA=\cA_1+\cA_2+ \cA_3$};
    \etik 
\end{center}

The professor then went back to what he was about to say, when Feynman chimed in again: "What if I now poke another hole in my screen? Then another, and another after that?" The professor at this point had had enough and simply said "Ok wise guy, I think it's clear to everyone that you simply sum over all the different path amplitudes. Now where was I..." But before he could continue, yup that's right Feynman chimed in again. "What about if I include another screen?" Slightly irritated, but bound by his profession to be polite, the professor drew the following diagram

\begin{center}
    \btik 
        \node at (0,0) {\Huge{$*$}};
        \node at (0,-0.5) {Source};
        \draw[thick] (2,2) -- (2,1.1);
        \draw[thick] (2,0.9) -- (2,-0.4);
        \draw[thick] (2,-0.6) -- (2,-0.9);
        \draw[thick] (2,-1.1) -- (2,-2);
        %
        \midarrow (0,0) -- (2,1);
        \midarrow (0,0) -- (2,-1);
        \midarrow (0,0) -- (2,-0.5);
        %
        \draw[thick] (4,2) -- (4,1.5);
        \draw[thick] (4,1.3) -- (4,-0.2);
        \draw[thick] (4,-0.4) -- (4,-1.3);
        \draw[thick] (4,-1.5) -- (4,-2);
        %
        \midarrow (4,1.4) -- (6,0.5);
        \midarrow (4,-0.3) -- (6,0.5);
        \midarrow (4,-1.4) -- (6,0.5);
        %
        \midarrow (2,1) -- (4,1.4);
        \midarrow (2,1) -- (4,-0.3);
        \midarrow (2,1) -- (4,-1.4);
        \midarrow (2,-1) -- (4,1.4);
        \midarrow (2,-1) -- (4,-0.3);
        \midarrow (2,-1) -- (4,-1.4);
        \midarrow (2,-0.5) -- (4,1.4);
        \midarrow (2,-0.5) -- (4,-0.3);
        \midarrow (2,-0.5) -- (4,-1.4);
        %
        \draw[fill=black] (6,0.5) circle [radius=0.07] node [right] {$\cA=\sum$All Paths};
    \etik 
\end{center}

Now before Feynman could chime in again the professor said "And before you ask, if I add another screen, I just do the same thing and sum over all paths again." However Feynman had one more, very important, chime in. "Ok Sir," he said, ignoring the frustrated look on the professor's face, "what if I include an infinite number of screens and on each screen I poke an infinite number of holes, so that in total there are no screens and no holes?" The professor just laughed and continued with the rest of his course. 

But as Zee points on, this last comment has a very important result. If we want to find the amplitude for a particle to propagate from some source and arrive at a detector, we must consider \textit{every single path} connecting the two, and sum over all their amplitudes.
\begin{center}
    \btik 
        \node at (0,0) {\Huge{$*$}};
        %\node at (0,-0.5) {Source};
        \draw[fill=black] (6,0.5) circle [radius=0.07];
        \midarrow (0,0) .. controls (2,1) and (4,-2) .. (6,0.5);
        \midarrow (0,0) .. controls (2,-3) and (4,-2) .. (6,0.5);
        \midarrow (0,0) .. controls (2,1) and (4,2) .. (6,0.5);
        \midarrow (0,0) .. controls (4,-1) and (5,3) .. (6,0.5);
        \midarrow (0,0) .. controls (3,-3) and (4,1) .. (6,0.5);
    \etik 
\end{center}
\noindent However there is obviously an infinite number of paths connecting the two and so the result is not a sum, but an integral. And thus the path integral is born.

\br 
\label{rem:SmoothPaths}
    On a technical note, we require the paths we integrate over to be smooth, that is continuous and infinitely differentiable. This is a reasonable thing to assume, as a discontinuous line would correspond to a particle `teleporting' while a non-continuous one would mean its velocity was discontinuous, which is also a problem.  
\er 

\subsection{Relation To Feynman Kernel}

Ok so we've answered the first question about what a path integral is, now we just need to see how it's related to the Feynman kernel. Well in the case of 1-dimensional QM the analogy is that each screen corresponds to a fixed time and the points each slit a position at that time. So the path between two slits is telling us the amplitude to go from one position at one time to another position at a later time. Well this is just the Feynman kernel. So in the `limit' that we remove all the screens, our Feynman kernel is an integral over all the different paths through spacetime connecting the initial position at initial time to the final position at the final time. 

\br 
    Note to go to a higher dimensional case is easy, simply replace the slits with a $d$-dimensional sheet with holes poked in it. 
\er 

\section{Hamiltonian Phase Space}

Let's calculate something now. It might not be obvious at first what we're calculating, but it will become clear next lecture. In order to do this, first you need to complete the following exercise.

\bbox 
    Given that in the position space representation, a wavefunction for a momentum eigenstate with momentum $p$ is 
    \bse 
        \psi_p(q) := \braket{q}{p}, \qquad -i\hbar\frac{d}{dt}\psi_p(q) = p \psi_p(q),
    \ese 
    show that 
    \be 
    \label{eqn:pqInnerProduct}
        \braket{p}{q} = (\braket{q}{p})^* = \frac{1}{\sqrt{2\pi\hbar}} \exp\bigg(-\frac{i}{\hbar} pq\bigg).
    \ee 
    \textit{Hint: Use the differential equation above to guess the form of $\psi_p(q)$ and then impose a normalisation condition.}
\ebox 

So what we want to find is 
\bse 
    \bra{p_F}U(T)\ket{q_I} = \bra{p_F} e^{-\frac{i}{\hbar}H}\ket{q_I}.
\ese 
Note this is \textit{not} the Feynman kernel as we have $p_F$ for the left hand state. With the above intuition in mind, let's break the the time interval up into $n$ equal intervals with 
\bse 
    t_i = t_I + \frac{\epsilon}{n}, \qquad \text{with} \qquad \epsilon = \frac{T}{n}.
\ese 
Note that this gives us $t_0=t_I$ and $t_n=t_F$, which is obviously what we want. We shall therefore relabel $p_F=p_n$ and $q_I=q_0$. Using the $\epsilon$ definition we can rewrite the our expression as
\bse 
    \bra{p_n} e^{(-\frac{i}{\hbar}H\epsilon)^n}\ket{q_0} = \bra{p_n} \underbrace{e^{-\frac{i}{\hbar}H\epsilon} e^{-\frac{i}{\hbar}H\epsilon} ... e^{-\frac{i}{\hbar}H\epsilon}}_{n\text{-times}} \ket{q_0} .
\ese 
Now we employ our trick of inserting the identity as a complete set of states. In fact we insert 
\bse 
    \b1 = \b1^2 = \int dq \ket{q}\bra{q} \int dp \ket{p}\bra{p} = \int \frac{dqdp}{\sqrt{2\pi\hbar}} e^{\frac{i}{\hbar}pq},
\ese 
where we have used \Cref{eqn:pqInnerProduct}. We do this in-between every exponential term. This gives us\footnote{And so begins the painful job of following terms around...}
\bse 
    \bra{p_n}U(T)\ket{q_0} = \int \Bigg[ \prod_{i=1}^{n-1} \frac{dq_idp_i}{\sqrt{2\pi\hbar}} e^{\frac{i}{\hbar}p_iq_i} \Bigg] \bra{p_n} e^{-\frac{i}{\hbar}H\epsilon}\ket{q_{n-1}} \bra{p_{n-1}} e^{-\frac{i}{\hbar}H\epsilon}\ket{q_{n-2}} ... \bra{p_1} e^{-\frac{i}{\hbar}H\epsilon}\ket{q_0}.
\ese 
We now take the limit $n\to\infty$ (an infinite number of screens), this gives $\epsilon\to 0$, and so we can Taylor expand the exponentials:
\bse 
    e^{-\frac{i}{\hbar}H\epsilon} = \b1 - \frac{i}{\hbar} H\epsilon + \cO\bigg(\frac{1}{n^2}\bigg).
\ese 
Then using \Cref{eqn:Hamiltonian}, we have\footnote{Here I reinsert the hats for clarity} 
\bse 
    \begin{split}
        \bra{p_i}\hat{H}\ket{q_{i-1}} & = \frac{1}{2m}\bra{p_i}\hat{q}^2\ket{q_{i-1}} + \bra{p_i} \hat{V}(\hat{q})\ket{q_{i-1}} \\
        & = \bigg(\frac{q_{i-1}^2}{2m} + V(q_{i-1})\bigg)\braket{p_i}{q_{i-1}} \\
        & = \frac{1}{\sqrt{2\pi\hbar}} H(p_i,q_{i-1}) e^{-\frac{i}{\hbar}p_iq_{i-1}}, 
    \end{split}
\ese 
which gives us 
\bse 
    \bra{p_i}e^{-\frac{i}{\hbar}\hat{H}\epsilon}\ket{q_{i-1}} = \frac{1}{\sqrt{2\pi\hbar}} \exp\Bigg(-\frac{i}{\hbar} \big( p_iq_{i-1} +\epsilon H(p_i,q_{i-1})\big)\bigg) + \cO\bigg(\frac{1}{n^2}\bigg).
\ese 
Putting this into our expression above, we get 
\bse 
    \begin{split}
        \bra{p_n}U(T)\ket{q_0} = \frac{1}{\sqrt{2\pi\hbar}} \int \Bigg[ \prod_{i=1}^{n-1}\frac{dq_idp_i}{2\pi\hbar}\bigg] \exp\Bigg\{ & \frac{i}{\hbar}\bigg[ \sum_{i=1}^{n-1} \bigg(\frac{p_i(q_i-q_{i-1})}{\epsilon} - H(p_i,q_{i-1})\bigg) \epsilon \\
        & - p_nq_{n-1} - H(p_n,q_{n-1})\epsilon\bigg]\Bigg\} + \cO\bigg(\frac{1}{n^2}\bigg). 
    \end{split}
\ese 
We then take the limit $n\to\infty$ and define the path integral measure
\mybox{
    \be 
    \label{eqn:PathIntegralMeasure}
        \int DpDq := \lim_{n\to\infty} \int \Bigg[ \prod_{i=1}^{n-1}\frac{dq_i}{\sqrt{2\pi\hbar}}\frac{dp_i}{\sqrt{2\pi\hbar}}\Bigg],
    \ee 
}
\noindent and then argue that $q_{n-1}=q_n$ in this limit and also use 
\bse 
    \lim_{n\to\infty} \frac{q_i-q_{i-1}}{\epsilon} = \dot{q}(t_i)
\ese 
i.e. the velocity of the particle,\footnote{Note this is well defined given \Cref{rem:SmoothPaths}} to finally give us
\mybox{
    \be 
    \label{eqn:pUq}
        \bra{p_n}U(T)\ket{q_0} = \frac{1}{\sqrt{2\pi\hbar}} \int DpDq \exp\bigg[ \frac{i}{\hbar} \bigg(\int_{t_I}^{t_F} dt \Big(p(t)\dot{q}(t) - H(p,q)\Big) - p(t_F)q(t_F)\bigg) \bigg].
    \ee
}

\section{Wick Rotation}

As promised above, next lecture we will see why we have bothered to derive \Cref{eqn:pUq}, but in order to do that we will need to use a very popular trick, the \textit{Wick rotation}, and we conclude this lecture with this. 

As we will see, we are going to take the limit $t\to\infty$ next lecture, and so we need to work out if such a thing is well defined for our time-evolution operator. That is we want to ask what happens to\footnote{Note we've inserted $\b1=\int dE \ket{E}\bra{E}$.}
\be 
\label{eqn:UWickRotation}
    U(t)\ket{\psi} = \int dE e^{-\frac{i}{\hbar}Ht} \ket{E}\braket{E}{\psi} = \int dE e^{-\frac{i}{\hbar}Et} \ket{E}\braket{E}{\psi}
\ee 
in the limit $t\to\infty$. Well consider the analytic continuation of $t$ to the complex plane:
\be 
\label{eqn:WickRotation}
    t \to e^{-i\theta} t, \qquad \theta\in [0,\pi/2].
\ee  
This is what we call a Wick rotation, and we note that $\theta=0$ changes nothing, but $\theta=\pi/2$ takes our exponential in \Cref{eqn:UWickRotation} from being complex to being real. People often refer to this as "going to Euclidean time", for obvious reasons. Ok so what happens now, well our expression becomes (setting $\hbar=1$ for this calculation)
\bse 
    U(t)\ket{\psi} = \int dE e^{-i\cos\theta Et} e^{-\sin\theta Et} \ket{E}\braket{E}{\psi}.
\ese 
Now if we assume we have a unique ground state $\ket{\Omega}$, scaled such that $\braket{\Omega}{\psi} \neq 0$ and that we have a mass gap,\footnote{That is the next energy level is strictly greater then $E_{\Omega}$} then the second exponential in the above expression suppresses all other energy states in our limit. That is, 
\be 
\label{eqn:UWickRotationLimit}
    \lim_{t\to\infty} U(t)\ket{\psi} = e^{-i\cos\theta E_{\Omega}t} e^{-\sin\theta E_{\Omega}t} \ket{\Omega}\braket{\Omega}{\psi}.
\ee 

\bbox 
    Convince yourself that similarly to \Cref{eqn:UWickRotationLimit}, we have 
    \bse 
        \lim_{t\to\infty} \bra{\psi}U(t) = e^{-i\cos\theta E_{\Omega}t} e^{-\sin\theta E_{\Omega}t} \bra{\Omega}\braket{\psi}{\Omega}.
    \ese 
\ebox 
\chapter{Green's Functions}

\section{$N$-Point Green's Function}

Last lecture we derived an expression for $\bra{p_n}U(T)\ket{q_0}$ and said we would use it for something. Let's now consider the same calculation, but now we insert a bunch of $q$ operators at times $t_1<t_2<...<t_N$ in between the different $U(t)$s. That is consider 
\bse 
    \bra{p_F} U(t_F-t_N) q \,  U(t_N-t_{N-1}) q  ... q \,U(t_1-t_I)\ket{q_I} = \bra{p_F} U(t_F) \cT\big(q(t_N)...q(t_1)\big) U(t_I)\ket{q_I}.
\ese 
This equation above says "start at some initial position and time, $\ket{q_I}$, then evolve to $t_1$, $U(t_1-t_I)$, then insert a $q$ operator at $t_1$, then evolve to $t_2$, $U(t_2-t_1)$, insert a $q$, etc until you get to $t_N$ and then finally evolve to $t_F$ and arrive at the final state." For clarity, what we're doing is considering all the paths from initial state $\ket{q_I}$ to final state $\ket{p_F}$ and fixing $N$ points along the way by the values $q(t_i)$. That is, every path not only has to start and end at the same place, but they must all meet at $N$ points in between. For example, if $N=2$ we have two fixed points, and all paths must meet there, as indicated with the two paths below. The dashed lines are our equal time sheets, and the two dots represent the insertion of an operator $q$. 
\begin{center}
    \btik 
        \draw[thick] (0,0) .. controls (2,1) and (4,-1) .. (6,0);
        \draw[thick] (0,0) .. controls (0.5,-0.5) and (2,1) .. (3.5,0) .. controls (4.5,-0.7) .. (6,0);
        \draw[fill=black] (1.67,0.27) circle [radius=0.07];
        \draw[fill=black] (3.75,-0.18) circle [radius=0.07];
        \draw[dashed] (0,1) -- (0,-1);
        \draw[dashed] (1.67,1) -- (1.67,-1);
        \draw[dashed] (3.75,1) -- (3.75,-1);
        \draw[dashed] (6,1) -- (6,-1);
        \node at (0,-1.2) {$t_I$};
        \node at (1.67,-1.2) {$t_1$};
        \node at (3.75,-1.2) {$t_2$};
        \node at (6,-1.2) {$t_F$};
    \etik 
\end{center}

The calculation will follow through exactly as above, but now the $q$ operators will act on the inserted states $\ket{q_i}$. This gives us 
\bse 
    \frac{1}{\sqrt{2\pi\hbar}} \int DpDq \, q(t_N)q(t_{N-1}) ... q(t_1) \exp\bigg[ \frac{i}{\hbar} \bigg(\int_{t_I}^{t_F} dt \Big(p(t)\dot{q}(t) - H(p,q)\Big) - p(t_F)q(t_F)\bigg) \bigg],
\ese 
where the $q(t_i)$s here are the eigenvalues.

\br 
\label{rem:PathIntegralTimeOrdered}
    Note the path integral automatically gives us the time-ordering. That is, even if we switch the $q(t_i)$s around in the final expression (which we can do because they're just numbers) we can arrive back at the time ordered expression. This is just because we define our evolution operator to only go forward in time, and so the only way we can order our time evolution operators is so that $q(t_N)$ appears furthest to the left and $q(t_1)$ furthest to the right.
\er 

We now use the result from the end of last lecture, to claim that in the limit $t_I\to-\infty$ and $t_F\to\infty$, the above simply becomes the so-called \textit{$N$-point Green's function}:
\bse 
    G(t_1,...,t_N) := \bra{\Omega} \cT \big(q(t_1)...q(t)N)\big) \ket{\Omega}.
\ese 
This just comes from taking the limit using the right hand side of the first equation in this lecture. So we get an explicit expression for it as 
\bse 
    G(t_1,...,t_N) = \widetilde{\cN} \int DpDq \, q(t_N)q(t_{N-1}) ... q(t_1) \exp \bigg[ \frac{i}{\hbar} \int_{-\infty e^{-i\theta}}^{{\infty e^{-i\theta}}} dt \, \big(p(t)\dot{q}(t) - H(p,q)\big)\bigg],
\ese 
for some normalisation constant $\widetilde{\cN}$, which contains the $p(t_F)q(t_F)$ term in the previous exponential. Note the limits on the integrals contain factors of $e^{-i\theta}$ to remind us that we're doing the Wick rotated integral. To save notation, we shall drop the integral limits below, but it's important to note that we have to take this limit to define the $N$-point Green's function.

\br 
\label{rem:EuclideanTime}
    With that last comment in mind, it is common to take $\theta=\pi/2$ so that $e^{-i\theta}=-i$ and then define $\tau=it$ so that we can write the integral as\footnote{Bonus exercise, check this.} 
    \bse 
        i\int_{i\infty}^{-i\infty} dt =  \int_{-\infty}^{\infty} d\tau
    \ese 
    As mentioned last lecture, $\tau$ is referred to as \textit{Euclidean time}. Note in Euclidean time $\dot{q}=i\dot{q}$, and so 
    \bse 
        L(q,\dot{q}) \to -\frac{m}{2}\dot{q}^2 - V(q),
    \ese 
    and so we define the \textit{Euclidean Action}
    \bse 
        S_E[q] := \int d\tau \bigg(\frac{m}{2}\dot{q}^2 + V(q)\bigg)
    \ese 
    This allows us to see the exponential suppression easier as we get (setting $\hbar=1$ for a moment). 
    \bse 
        e^{iS} \to e^{-S_E}.
    \ese 
\er 

\subsection{1-Dimensional QM}

Let's evaluate this for our 1-dimensional QM system. We have
\bse 
    \begin{split}
        p\dot{q} - H & = p\dot{q} - \frac{p^2}{2m} - V(q) \\
        & = -\frac{1}{2m}(p-m\dot{q})^2 + \frac{m}{2}\dot{q}^2 - V(q) \\
        & = - \frac{1}{2m} (p-m\dot{q})^2 + L(q,\dot{q}),
    \end{split}
\ese 
which allows us to split the $Dp$ and $Dq$ integral, giving us 
\bse 
    G(t_1,...,t_N) = \widetilde{N} \int Dp \exp\bigg( -\frac{i}{2m\hbar}\int dt \, (p-m\dot{q})^2 \bigg) \int Dq \, q(t_N)...q(t_1) \exp \bigg[ \frac{i}{\hbar} \int dt \, L(q,\dot{q}) \bigg].
\ese

\bcl 
    The integral 
    \bse 
        \int Dp \exp\bigg( -\frac{i}{2m\hbar}\int dt \, (p-m\dot{q})^2 \bigg)
    \ese 
    is well defined. 
\ecl 
\bq 
    Given \Cref{rem:EuclideanTime}, at first site this claim doesn't seem true, as going to Euclidean time just removes the $i$ and puts a minus sign in, which in this case would give us a positive exponential. However we have to remember that $\dot{q}$ involves a time derivative and so is also effected. We therefore prove the claim by considering discreting time. The integral then becomes a product of integrals of the form
    \bse 
        \int d\widetilde{p} \exp\bigg[ -\frac{i\epsilon}{2m\hbar} \widetilde{p}^2\bigg],
    \ese 
    where we have defined 
    \bse 
        \widetilde{p} := (p-m\dot{q}),
    \ese 
    and used the fact that $p$ and $\widetilde{p}$ are linearly related to change the measure to $d\widetilde{p}$. Under the Wick rotation, this becomes 
    \bse 
        \int d\widetilde{p} \exp\bigg[ -\frac{\epsilon}{2m\hbar}\big(\sin\theta +i\cos\theta\big) \widetilde{p}^2\bigg] = \sqrt{\frac{2\pi\hbar m}{i\epsilon(\cos\theta-i\sin\theta )}},
    \ese
    where the right-hand side comes from Gaussian integral formula which states 
    \bse 
        \int dx e^{-ax^2} = \sqrt{\frac{\pi}{a}} \qquad \text{if } \qquad \Re a >0,
    \ese 
    which we have as $\theta \in [0,\pi/2]$. If we then set $\theta=0$, we get 
    \be 
    \label{eqn:DpGaussianResult}
        \int Dp \exp\bigg(-\frac{i}{2m\hbar} \int dt(p-m\dot{q})^2\bigg) = \lim_{n\to\infty} \bigg(\frac{-im}{\epsilon}\bigg)^{\frac{n-1}{2}},
    \ee 
    where have used the definition of $Dp$, \Cref{eqn:PathIntegralMeasure}.
\eq 

So if we absorb this result into our normalisation constant, and use the definition of the action, we get 
\bse 
    G(t_1,...,t_N) = \cN \int Dq \, q(t_N)...q(t_1) \exp\bigg(\frac{i}{\hbar} S[q]\bigg).
\ese 

We can find the normalisation constant by consider the $0$-point Green's function, which is just the vacuum inner product $\braket{\Omega}{\Omega} =1$, as we define it to be normalised. Therefore we conclude 
\bse
    \cN = \frac{1}{\int Dq \, \exp\Big(\frac{i}{\hbar} S[q]\Big)}, 
\ese 
finally giving us the result 
\mybox{
\be 
\label{eqn:NPointGreensFunctionQM}
    G(t_1,...,t_N) = \frac{\int Dq \, q(t_N)...q(t_1) \exp\Big(\frac{i}{\hbar} S[q]\Big)}{\int Dq \, \exp\Big(\frac{i}{\hbar} S[q]\Big)}.
\ee 
}

\br 
    Note that $G(t_1,...,t_N)$ is totally symmetric in all its entries, as the $q(t_i)$s are just numbers. Recall \Cref{rem:PathIntegralTimeOrdered} tells us we can do this without altering the time ordering we need.  
\er 

In Euclidean time this can be written 
\bse 
    G_E(t_1,...,t_N) = \frac{\int Dq \, q(t_N)...q(t_1) \exp\Big(-\frac{1}{\hbar} S_E[q]\Big)}{\int Dq \, \exp\Big(-\frac{1}{\hbar} S_E[q]\Big)}
\ese

\section{Partition Function}

\subsection{Functional Derivative}

Recall that a functional is an object that maps functions to numbers. The action is an example, as it maps the functions $q(t)$ to some number. Recall the derivative of a function is given by 
\bse 
    \frac{df}{dx} := \lim_{\epsilon\to0} \frac{f(x+\epsilon)-f(x)}{\epsilon}.
\ese
We can use this to motivate the definition of the derivative of a functional. 

\bd[Functional Derivative]
    Given a functional $F[J]$, we define the \textit{functional derivative} of $F$ as 
    \be 
    \label{eqn:FunctionalDerivative}
        \frac{\del F[J]}{\del J(s)} := \lim_{\epsilon\to0} \frac{F[J(t)+\epsilon\del(t-s)]-F[J]}{\epsilon}.
    \ee 
\ed 

We are often only interested in the variation $\del F[J]$, which is given by
\be 
\label{eqn:VaritaionOfFunctional}
    \del F[J] = \int ds \, \del J(s) \, \frac{\del F[J]}{\del J(s)}.
\ee 
We can also find the variation by simply calculating 
\bse 
    F[J+\del J] - F[J]
\ese 
to order $\del J$.

\bc 
    It follows immediately from the definition above that for the functional 
    \bse 
        F[J] = J(t), \qquad \forall J
    \ese 
    where $t$ is some fixed value in the argument of $J$ that 
    \bse 
        \frac{\del J(t)}{\del J(s)} := \frac{\del F[J]}{\del J(s)} = \del(t-s),
    \ese 
    where the left-hand expression is defined for notational simplicity.
\ec 

\br 
    We can treat $\frac{\del}{\del J(s)}$ sort of like a regular partial derivative, and use things like the product rule. However we should be careful when doing this as it doesn't always work.
\er 

\bex 
    Let 
    \bse 
        F[J] = \int dt \phi(t) J(t),
    \ese 
    then
    \bse 
        \begin{split}
            \frac{\del F[J]}{\del J(s)} & = \lim_{\epsilon\to0} \frac{\int dt \phi(t) \big(J(t) + \epsilon\del(t-s)\big) - \int dt \phi(t) J(t)}{\epsilon} \\
            & = \phi(s)
        \end{split}
    \ese 
    We get the same result using the partial derivative approach:
    \bse 
        \begin{split}
            \frac{\del F[J]}{\del J(s)} & = \int dt \bigg(\frac{\del \phi(t)}{\del J(s)}J(t) + \phi(t) \frac{\del J(t)}{\del J(s)}\bigg) \\
            & = \int dt \phi(t) \del(t-s) \\
            & = \phi(s).
        \end{split}
    \ese
    This gives us 
    \bse 
        \del F[J] = \int ds \phi(s) \del J(s).
    \ese 
\eex

\subsection{Partition Function}

We can take the sum over all the $N$-point Green's functions and get what is known as a \textit{generating functional}. For the specific case of the sum of $N$-point Green's functions, we call the resulting generating functional the \textit{partition function}.
\mybox{
    \be 
    \label{eqn:PartitionFunction}
        Z[J] := 1 + \sum_{N=1}^{\infty} \frac{1}{N!} \int ds_1... ds_N G(s_1,...,s_N) J(s_1) ... J(s_N).
    \ee 
}
\noindent The partition function is incredible useful as it allows us to calculate \textit{any} $N$-point Green's function by taking $N$ functional derivatives and then setting $J=0$. That is 
\mybox{
    \be
    \label{eqn:GreensFunctionsFromParition}
        G(t_1,..,t_N) = \frac{\del^N Z[J]}{\del J(t_1) ... \del J(t_N)} \bigg|_{J=0}.
    \ee 
}

\bex 
    Let's do it for $N=2$. The derivatives will kill the first two terms in the expansion, and then anything with more than two $J$s will vanish when we set $J=0$. So we are left with 
    \bse 
        \begin{split}
            \frac{\del^2Z[J]}{\del J(t_1) \del J(t_2)}\bigg|_{J=0} & = \frac{1}{2} \int ds_1 ds_2 G(s_1,s_2) \frac{\del}{\del J(s_1)} \big( \del(s_1-t_2) J(s_2) + J(s_1) \del(s_2-t_2)\big) \\
            & = \frac{1}{2}\int ds_1ds_2 G(s_1,s_2) \big(\del(s_1-t_2)\del(s_2-t_2) + \del(s_2-t_2)\del(s_1-t_2)\big) \\
            & = \frac{1}{2}\big(G(t_2,t_1) + G(t_1,t_2)\big) \\
            & = G(t_1,t_2),
        \end{split}
    \ese 
    where the last line has used the symmetric property of the $N$-point Green's functions.
\eex 

\subsection{Path Integral Form}

We can use the path integral expression for the $N$-point Green's function, \Cref{eqn:NPointGreensFunctionQM}, to write the partition function as a path integral. We have 
\bse 
    \begin{split}
        \frac{1}{\cN} Z[J] & = \int Dq \bigg( e^{\frac{i}{\hbar}S[q]} + \sum_{N=1}^{\infty} \frac{1}{N!} \int ds_1...ds_N J(s_1) ... J(s_N) q(s_1) ... q(s_N) \,  e^{\frac{i}{\hbar}S[q]}\bigg) \\
        & = \int Dq \sum_{N=0}^{\infty} \frac{1}{N!} \bigg(\int ds \,  q(s)J(s)\bigg)^N e^{\frac{i}{\hbar}S[q]} \\
        & = \int Dq \exp \bigg( \frac{i}{\hbar}S[q] + \int ds \, q(s) J(s)\bigg),
    \end{split}
\ese 
so that 
\mybox{
    \be
    \label{eqn:PartitionFunctionPathIntegral}
        Z[J] = \cN \int Dq \exp \bigg( \frac{i}{\hbar}S[q] + \int dt \, q(t) J(t)\bigg)
    \ee
}

\br 
    \Cref{eqn:PartitionFunctionPathIntegral} is often defined using another convention given by sending
    \bse 
        \frac{\del}{\del J(t)} \to -i\hbar \frac{\del}{\del J(t)},
    \ese
    so that we get 
    \bse 
        Z[J] = \cN \int Dq \exp \bigg[ \frac{i}{\hbar}\bigg( S[q] + \int dt \, q(t) J(t)\bigg)\bigg].
    \ese 
    We can also write it as a Euclidean path integral. If we use the convention \Cref{eqn:PartitionFunctionPathIntegral}, we redefine $J(t)\to iJ(t)$ to give 
    \bse 
        Z[J] = \cN \int Dq \exp \bigg( -\frac{1}{\hbar}S_E[q] + \int dt \, q(t) J(t)\bigg).
    \ese 
\er 

\section{Back To The Feynman Kernel}

Ok now let's return to the Feynman kernel and express it as a path integral. Recall that it's defined as 
\bse 
    K(q_F,T;q_I,0) := \bra{q_F}U(T)\ket{q_I}.
\ese 
We can insert the complete set of states
\bse 
    \b1 = \int dp_F \ket{p_F}\bra{p_F} 
\ese 
to give
\bse 
    \begin{split}
        K(q_F,T;q_I,0) & = \int dp_F \braket{q_F}{p_F}\bra{p_F}U(T)\ket{q_I} \\
        & =  \lim_{n\to\infty} \bigg(\frac{-im}{\epsilon}\bigg)^{\frac{n-1}{2}} \int dp_F \frac{1}{\sqrt{2\pi\hbar}}e^{\frac{i}{\hbar}p_Fq_F} \frac{1}{\sqrt{2\pi\hbar}}\int Dq \exp\bigg(\frac{i}{\hbar} S[q] - p_Fq_{n-1}\bigg),
    \end{split}
\ese 
where we have used \Cref{eqn:pqInnerProduct,eqn:DpGaussianResult,eqn:pUq}, with the latter manipulated as we did for the Green's function to get the action in the exponential. Note that we haven't taken the limit on the $q_{n-1}$ in the exponential yet, this is because we can now combine it with the first exponential and use 
\bse 
    \int dp_F \exp\bigg( \frac{i}{\hbar} p_F(q_F-q_{n-1})\bigg) = (2\pi\hbar) \del(q_F-q_{n-1}),
\ese
to give (putting the definition of $Dq$, \Cref{eqn:PathIntegralMeasure}, in)
\be 
\label{eqn:FeynmanKernerlDiscreteTime}
    \begin{split}
        K(q_F,T;q_I,0) & = \lim_{n\to\infty}\bigg(\frac{-im}{\epsilon}\bigg)^{\frac{n-1}{2}}  \int \Bigg[ \prod_{i=1}^{n-1}\frac{dq_i}{\sqrt{2\pi\hbar}}\bigg] \del(q_F-q_{n-1}) \exp\bigg(\frac{i}{\hbar} S[q]\bigg) \\
        & =  \lim_{n\to\infty} \sqrt{\frac{-im}{2\pi\hbar \epsilon}}  \int\Bigg[\prod_{i=1}^{n-2}\sqrt{\frac{-im}{2\pi\hbar \epsilon}} dq_i\Bigg] \exp\bigg(\frac{i}{\hbar} S[q]\bigg)\bigg|_{q_{n-1}=q_F},
    \end{split}
\ee
where on the second line we have done some rearrangement (notice the upper limit on the product changes because of delta function). So if we redefine $Dq$ appropriately, we get 
\mybox{
    \be
    \label{eqn:FeynmanKernelPathIntegral}
        K(q_F,T;q_I,0) = \int Dq \exp\bigg(\frac{i}{\hbar}S[q]\bigg)
    \ee 
}
\chapter{Classical Limit \& Schwinger-Dyson Equation}

\section{Classical Limit}

Last lecture by showing that in our path integral, every path contributes a factor of $e^{\frac{i}{\hbar}S[q]}$, all with magnitude $1$. This was a result in the quantum theory, and as always it is instructive to check that in some limit we get the classical theory back, i.e. only paths that satisfy the Euler-Lagrange equations. We shall do just that now. 

When we observe a path, there is always some uncertainty about the path, that is we cannot distinguish between $q(t)$ and $q(t)+\eta(t)$, where $\eta(t)$ is some small deviation. Now the action for this latter path is given by 
\begin{equation*}
    \begin{split}
        S[q+\eta] & = \int_{t_I}^{t_F} dt \, L(q+\eta, \dot{q}+\dot{\eta}) \\
        & = \int_{t_I}^{t_F} dt \Bigg( L(q,\dot{q}) + \eta(t) \frac{\p L}{\p q}(q,\dot{q}) + \dot{\eta}(t) \frac{\p L}{\p \dot{q}}(q,\dot{q}) + \cO(\eta^2)\bigg) \\
        & = S[q] + \int_{t_I}^{t_F} dt \, \eta(t) \bigg( \frac{\p L}{\p q} -\frac{d}{dt} \bigg[\frac{\p L}{\p \dot{q}}\bigg]\bigg) + \cO(\eta^2),
    \end{split}
\end{equation*}
where we've used the standard tool of integrating by parts and saying that $\eta(t_I)=0=\eta(t_F)$, and where on the last line we have suppressed the arguments of $L$ for notational reasons. So we see the deviated path contributes a phase factor of 
\bse 
    \int_{t_I}^{t_F} dt \, \eta(t) \bigg( \frac{\p L}{\p q} -\frac{d}{dt} \bigg[\frac{\p L}{\p \dot{q}}\bigg]\bigg) + \cO(\eta^2).
\ese 
So what does this mean for the classical limit? Well we need to sum over all the different $\eta(t)$ variations, and if $q(t)$ is not the classical path $q_{cl}(t)$, i.e. the one that satisfies the Euler-Lagrange equations, we get a wide range of contributions to the phase. Overall these will approximately destructively interfere and so will not contribute to drastically to the result. Now if $q(t)=q_{cl}(t)$ then the bracket term vanishes and the contributions to the phase are $\cO(\eta^2)$, and so the range is decreased. We claim that this results in a significant reduction in the destructive interference and contributes to constructive interference, and so this gives the biggest contribution. In other words, paths that aren't the classical path cancel whereas the ones near the classical path add, so we get the classical result. 

\br 
    We can also see this result by doing our Wick rotation to the Euclidean picture. Here we get contributions of the form 
    \bse 
        \exp\bigg( -\frac{1}{\hbar}\big(S_E[q]- S_E[q_{cl}]\big)\bigg),
    \ese 
    and so we see straight away that non-classical paths are exponentially suppressed.
\er 

\subsection{Variations Around The Classical Path}

What happens to our path integral when we take a variation around the classical path? Well just seen how the action changes, and showed that if we take the variation around the classical path we get 
\bse 
    S[q+\eta] = S[q_{cl}] + \cO(\eta^2).
\ese 
What about the integral measure $Dq$? Well to this the change $q=q_{cl}+\eta$ is just a change of variables, and we take $q_{cl}$ to be a `constant' so we just get 
\bse 
    Dq \to D\eta. 
\ese 
So the path integral is 
\bse 
    \int Dq e^{\frac{i}{\hbar}S[q]} = e^{\frac{i}{\hbar}S[q_{cl}]} \int D\eta \, e^{\frac{i}{\hbar}\cO(\eta^2)}.
\ese 
Now we note that if we consider an action that is at most quadratic in $q$, then our order $\eta^2$ terms are independent of $q_{cl}$ (we would need at least a cubic $q$ for that), and so for these cases we get 
\bse 
    \int Dq e^{\frac{i}{\hbar}S[q]} = e^{\frac{i}{\hbar}S[q_{cl}]} \int D\eta \, e^{\frac{i}{\hbar}S[\eta]}.
\ese 
We then simplify notation by using $t_I=0$, $t_F=T$ which gives
\be 
\label{eqn:eta(t)Sin}
    \eta(0)=0=\eta(T) \qquad \implies \qquad \eta(t) = \sum_{i=1}^{\infty} a_{\l} \sin\bigg(\frac{\l\pi t}{T}\bigg),
\ee 
then if we discretise the time over $n$ intermediate paths $q_j$, this is equivalent (up to a normalisation factor) to taking $\eta(t)$ as above and integrating over the $a_{\l}$. We shall return to this soon when consdiering the harmonic oscillator.

\section{Calculating The PI Exactly}

As anyone familiar with QFT knows, it is only in very special cases that we can obtain exact results and in all other cases we tend to turn to perturbation theory. We will address this next lecture, but for this lecture we will present particular examples of when we can obtain an exact result for the path integral. It is worth emphasising again that this is \textit{not} something you can do in general, but is a particular property of the systems we're studying. 

\subsection{Free Particle}

Of course the most simple system we can study is the free particle. This corresponds to setting $V(q)=0$, and so our Feynman kernel (in the discrete time formalism, i.e. using \Cref{eqn:FeynmanKernerlDiscreteTime}) becomes
\be 
\label{eqn:FeynmanKernelFreeParticleOne}
    K(q_F,T;q_I,0) = \lim_{n\to\infty} \sqrt{\frac{-im}{2\pi\hbar \epsilon}}  \int\Bigg[\prod_{i=1}^{n-2}\sqrt{\frac{-im}{2\pi\hbar \epsilon}} dq_i\Bigg] \exp\bigg(\frac{im}{2\hbar\epsilon} \sum_{j=1}^{n-1} (q_j-q_{j-1})^2 \bigg),
\ee 
where we also impose $q_0=q_I$ and $q_{n-1}=q_F$. We now notice that these are essentially a (huge) collection of Gaussian integrals, and so we can solve it. However, we need to be a bit clever because each $q_j$ appears twice, i.e. once as $q_j$ and once as $q_{(j+1)-1}$. We therefore want some kind of inductive proof of a result.

\bcl 
    We can rewrite \Cref{eqn:FeynmanKernelFreeParticleOne} as 
    \bse 
        K(q_F,T;q_I,0) = \lim_{n\to\infty} \sqrt{\frac{-im}{2\pi\hbar\epsilon}} \Bigg[ \prod_{j=1}^{n-2} \sqrt{\frac{j}{j+1}} \Bigg] \exp \bigg( \frac{im}{2\hbar\epsilon} \frac{1}{n-1} (q_F-q_I)^2\bigg)
    \ese 
\ecl 

\bq 
    First we note that the we can do the $dq_j$ integral without effecting any of the $k>j$ results but it will effect the $\ell<j$ ones. So we start at $j=1$ and work upwards. The proof then follows by noticing that we can write the $j=1$ term as (note $j_0=j_I$).
    \bse 
        \cI_1 = \sqrt{\frac{-im}{2\pi\hbar\epsilon}} \int dq_j \exp \bigg[ \frac{im}{2\hbar\epsilon} \bigg( (q_{j+1}-q_j)^2 + \frac{1}{j}(q_j-q_I)^2\bigg)\bigg].
    \ese 
    Let's see what this evaluates to. We start by expanding 
    \bse 
        (q_{j+1}-q_j)^2 + \frac{1}{j}(q_j-q_I)^2 = \bigg(\frac{j+1}{j}\bigg) q_j^2 -2\bigg(q_{j+1} + \frac{1}{j}q_I\bigg)q_j + q_{j+1}^2 + \frac{1}{j}q_I^2.
    \ese 
    Now we use a Wick rotation $T\to -iT$, and hence $\epsilon = T/n \to -i\epsilon$, so our integral becomes 
    \bse 
        \cI_1^E = \sqrt{\frac{m}{2\pi\hbar\epsilon}} \int dq_j \exp \big( -aq_j^2 + bq_j + c\big)
    \ese 
    where 
    \bse 
        a = \frac{m}{2\hbar\epsilon}\bigg(\frac{j+1}{j}\bigg), \qquad b = \frac{m}{\hbar\epsilon}\bigg(q_{j+1}+ \frac{1}{j}q_I\bigg), \qand c = -\frac{m}{2\hbar\epsilon}\bigg(q_{j+1}^2 + \frac{1}{j}q_I^2\bigg).
    \ese
    Now use the Gaussian result 
    \bse 
        \int dx \exp\big(-ax^2+bx+c\big) = \sqrt{\frac{\pi}{a}} \exp\bigg(\frac{b^2}{4a}+c\bigg), \qquad \text{if} \qquad \Re a >0,
    \ese 
    to give 
    \bse
        \cI_1^E = \sqrt{\frac{j}{j+1}} \exp\Bigg[ \frac{m}{2\hbar\epsilon}\Bigg( \frac{j}{j+1} \bigg[q_{j+1} + \frac{1}{j}q_I\bigg]^2 - \bigg[ q_{j+1}^2 + \frac{1}{j}q_I^2 \bigg] \Bigg) \Bigg].
    \ese
    Let's now focus on the bit inside round brackets: 
    \bse 
        \begin{split}
            \frac{j}{j+1} \bigg(q_{j+1} + \frac{1}{j}q_I\bigg)^2 - \bigg( q_{j+1}^2 + \frac{1}{j}q_I^2 \bigg) & = q_{j+1}^2\bigg(\frac{j}{j+1}-1\bigg) + \frac{q_I^2}{j}\bigg(\frac{1}{j+1}-1\bigg) - \frac{2q_{j+1}q_I}{(j+1)} \\
            & = -\frac{1}{j+1} \big( q_{j+1}^2 + q_I^2 - 2q_{j+1}q_I\big) \\
            & = -\frac{1}{j+1} (q_{j+1}-q_I)^2,
        \end{split}
    \ese 
    so if we undo our Wick rotation, i.e. $\epsilon\to i\epsilon$, we get 
    \bse 
        \cI_1 = \sqrt{\frac{j}{j+1}} \exp \bigg( \frac{im}{2\hbar\epsilon}\frac{1}{j+1} (q_{j+1}-q_I)^2\bigg).
    \ese
    
    So now we consider the $d1_2$ integral: the square root factor obviously just factors out and the rest is simply
    \bse 
        \exp\bigg[\frac{im}{2\hbar\epsilon} \bigg( \big(q_{j_2+1} - q_{j_2}\big)^2 + \frac{1}{j_1+1}\big(q_{j_1+1}-q_I
        \big)^2\bigg)\bigg],
    \ese 
    then we use $j_2=j_1+1$ to give 
    \bse 
        \cI_2 = \sqrt{\frac{j_1}{j_2}} \sqrt{\frac{-im}{2\pi\hbar\epsilon}} \int dq_2 \exp\bigg[\frac{im}{2\hbar\epsilon} \bigg( \big(q_{j_2+1} - q_{j_2}\big)^2 + \frac{1}{j_2}\big(q_{j_2}-q_I
        \big)^2\bigg)\bigg],
    \ese 
    which is (up to the first root factor) the equation we started the proof of. So we see the formula holds inductively (that's the reason we've been using $j$ all the time rather then setting $j=1$). 
    
    So we will just get a $\sqrt{j/j+1}$ factor for each integral, there's $n-2$ of these giving us the full result 
    \bse 
        K(q_F,T;q_I,0) = \lim_{n\to\infty} \sqrt{\frac{-im}{2\pi\hbar\epsilon}} \Bigg[ \prod_{j=1}^{n-2} \sqrt{\frac{j}{j+1}} \Bigg] \exp \bigg( \frac{im}{2\hbar\epsilon} \frac{1}{n-1} (q_{j-1}-q_I)^2\bigg),
    \ese 
    then we finally use $q_{j-1}=q_F$ which gives us exactly the result in the claim. 
\eq 

Now we note that the product in the above formula will simply give 
\bse 
    \prod_{j=1}^{n-2}\sqrt{\frac{j}{j+1}} = \frac{1}{\sqrt{n-1}},
\ese
and so, using $T = n\epsilon \sim  (n-1)\epsilon$, where the second term is understood in the limit $n\to\infty$, we finally get 
\be 
\label{eqn:FeynmanFernalFreeParticle}
    K(q_F,T;q_I,0) = \sqrt{\frac{-im}{2\pi\hbar T}} \exp\bigg( \frac{im}{2\hbar T} (q_F-q_I)^2\bigg) \propto \exp\bigg(\frac{i}{\hbar} S[q_{cl}]\bigg),
\ee 
where the second part comes from the fact that the square-root term is independent of $q_F/q_I$ and the definition of the classical action. 

\subsection{Simple Harmonic Oscillator}

Another case we can solve exactly is the simple harmonic oscillator. The Lagrangian is 
\bse 
    L = \frac{m}{2} \dot{q}^2 - \frac{m}{2}\omega^2q^2. 
\ese
This has classical action 
\be
\label{eqn:SHOClassicAction}
    S[q_{cl}] = \frac{m\omega}{2\sin(\omega T)}\Big[ (q_F^2+q_I^2)\cos(\omega T) -2q_Fq_I \Big]
\ee 

We can treat this as a variation to the classical path, and so as per the section above our path integral is
\bse 
    \begin{split}
        \int Dq e^{\frac{i}{\hbar}S[q]} & = e^{\frac{i}{\hbar}S[q_{cl}]} \int D\eta \, \exp\bigg( \frac{im}{2\hbar}\big( \dot{\eta}^2 - \omega^2\eta^2\big)\bigg) \\
        & = e^{\frac{i}{\hbar}S[q_{cl}]} \int D\eta \, \exp\bigg[ \frac{im}{2\hbar}\eta\bigg( -\frac{d^2}{dt^2} - \omega^2\bigg)\eta\bigg],
    \end{split}
\ese 
where the second line follows from integration by parts. 

\bbox
    Prove that integration by parts leads to the expression above.
\ebox  

We then use our simplified notation \Cref{eqn:eta(t)Sin} and note that this is an eigenvector of the operator in our integral, i.e. 
\bse 
    \bigg( -\frac{d^2}{dt^2} - \omega^2\bigg)\eta = \bigg(\frac{\l^2\pi^2}{T^2} - \omega^2\bigg) \eta,
\ese 
to obtain (absorbing everything else into the proportionality constant, we will get it all back at the end)
\bse
    \int Dq e^{\frac{i}{\hbar}S[q]} \propto \lim_{n\to\infty} \bigg[\prod_{j=1}^{n} \int da_{\l}\bigg] \exp\bigg[ \frac{im}{2\hbar}\bigg(\frac{\l^2\pi^2}{T^2} -\omega^2\bigg)\sum_{i=1}^n a^2_{\l}\bigg].
\ese 
Then we notice this is just the product of a bunch of Gaussian integrals, so we get
\bse
    \begin{split}
        \int Dq \, e^{\frac{i}{\hbar}S[q]} \propto \prod_{\l=1}^{\infty} \bigg(\frac{\l^2\pi^2}{T^2} -\omega^2\bigg)^{-1/2} \propto \prod_{\l=1}^{\infty} \bigg(1- \frac{\omega^2T^2}{\l^2\pi^2}\bigg)^{-1/2} = \bigg(\frac{\sin(\omega T)}{\omega T}\bigg)^{-1/2}.
    \end{split}
\ese 

Now the proportionality constant appears to have got very complicated as we've gone along, however it can easily be checked that we haven't removed anything that is a function of $\omega$, and so we know that our path integral is of the form 
\bse 
    \int D\eta \, e^{\frac{i}{\hbar}S[\eta]} = f(T) \bigg(\frac{\sin(\omega T)}{\omega T}\bigg)^{-1/2}.
\ese 
So how do we find $f(T)$, well we use 
\bse 
    \lim_{\omega\to0} \bigg(\frac{\sin(\omega T)}{\omega T}\bigg) = 1,
\ese 
to notice that in this limit our classic action, \Cref{eqn:SHOClassicAction}, is just that of a free particle action, which gives us
\bse 
    \lim_{\omega\to0} \int Dq \, e^{\frac{i}{\hbar}S[q]} = f(T) \exp\bigg(\frac{im}{2\hbar T} (q_F-q_I)^2 \bigg),
\ese 
which comparing to \Cref{eqn:FeynmanFernalFreeParticle} gives us 
\bse 
    f(T) = \sqrt{\frac{-im}{2\pi\hbar T}},
\ese
and finally
\bse 
    \int Dq \, e^{\frac{i}{\hbar}S[q]} = \sqrt{\frac{-im\omega}{2\pi\hbar T \sin(\omega T)}} \exp\bigg(\frac{m\omega}{2\sin(\omega T)}\Big[ (q_F^2+q_I^2)\cos(\omega T) -2q_Fq_I \Big]\bigg).
\ese 

\section{Schwinger-Dyson Equation}

As we said before, it is not normally true that we can solve path integrals exactly, and we need some method to solve them perturbatively. We can gain an appreciation of this fact by finding the differential equations that the partition function satisfies. Recall that the partition function can be written 
\bse 
    Z[J] = \cN \int Dq \exp \bigg( \frac{i}{\hbar} S[q] + \int dt q(t) J(t)\bigg).
\ese
In this expression, $q(t)$ is just an integration variable, and so the result is completely unchanged if we change variables to 
\bse 
    q(t) \to q(t) + \eta(t),
\ese 
provided $\eta(t)$ is independent of $q$. That is, we also have (using $D(q+\eta) = Dq$ for fixed $\eta$)
\bse 
    \begin{split}
        Z[J] & = \cN \int Dq \exp \bigg( \frac{i}{\hbar} S[q+\eta] +  \int dt q(t) J(t) + \int dt \eta(t) J(t)\bigg) \\
        & = \cN \int Dq \exp \bigg( \frac{i}{\hbar} S[q] + \frac{i}{\hbar}\int dt \eta(t) \frac{\del S[q]}{\del q(t)} +  \int dt q(t) J(t) + \int dt \, \eta(t) \frac{\del}{\del q(t)} \int ds J(s)q(s)\bigg) + \cO(\eta^2),
    \end{split}
\ese 
where the second line simply comes from the expanding $S[q+\eta]$ and rewriting 
\bse 
    \int dt \, \eta(t) J(t) = \int dt \, \eta(t) \frac{\del}{\del q(t)} \int ds J(s)q(s).
\ese 
Now the difference between the two different expressions for $Z[J]$ must vanish (as they are equal, all we've done is change integration variable). So if we take them away from each other and Taylor expand we get 
\bse 
    \begin{split}
        0 & =\frac{i}{\hbar} \int dt \, \eta(t) \frac{\del S[q]}{\del q(t)} + \int dt \, \eta(t) \frac{\del}{\del q(t)} \int ds \, J(s) q(s) + \cO(\eta^2) \\
        & = \int dt \eta(t) \frac{\del}{\del q(t)} \bigg( \frac{i}{\hbar}S[q] + \int ds J(s) q(s) \bigg) + \cO(\eta^2) \\
        & = \cN\int Dq \int dt \eta(t) \frac{\del}{\del q(t)} \exp\bigg( \frac{i}{\hbar}S[q] + \int ds J(s) q(s) \bigg) + \cO(\eta^2),
    \end{split}
\ese 
where the last line follows from the fact that the exponential is non-degenerate and that the path integral of 0 is 0. Now we use that we haven't specified the form of $\eta(t)$, and so we can conclude 
\bse 
    0 = \cN\int Dq \frac{\del}{\del q(t)} \exp\bigg( \frac{i}{\hbar}S[q] + \int ds J(s) q(s) \bigg).
\ese
Now we shall assume our Lagrangians are of the form 
\bse 
    L(q,\dot{q}) = \frac{m}{2}\dot{q}^2 + V(q),
\ese 
where $V(q)$ is some polynomial function. Then using the fact that the functional variation of the action gives us the Euler-Lagrange equations, the above condition become 
\bse 
    0 = \cN\int Dq \bigg( \frac{i}{\hbar} \Big[-V'\big(q(t)\big) - m \Ddot{q}(t) \Big] + J(t) \bigg) \exp\bigg( \frac{i}{\hbar}S[q] + \int ds J(s) q(s) \bigg).
\ese 
Finally we notice that 
\bse 
    \frac{\del}{\del J(t)}\exp\bigg( \frac{i}{\hbar}S[q] + \int ds J(s) q(s) \bigg) = q(t),
\ese 
and so we can replace the $q$s in the previous equation with functional derivatives w.r.t. $J(t)$. This allows us to pull these terms outside the path integral (i.e. past the $Dq$), giving us a differential equation for the partition function, known as the \textit{Schwinger-Dyson equation}: 
\mybox{
    \be
    \label{eqn:SchwingerDyson}
        \Bigg( -\frac{i}{\hbar} \Bigg[V'\bigg(\frac{\del}{\del J(t)}\bigg) + m \frac{d^2}{dt^2}\frac{\del}{\del J(t)} \Bigg] + J(t) \Bigg) Z[J] = 0.
    \ee 
}

This now explains to us why the previous two problems were solvable: if we have a quadratic action, then $V'$ is linear and so we get $2^{\text{nd}}$ order linear functional differential equation, which we can solve. However for the more general cases, when $V'$ is not linear, things get a lot more complicated. 
\chapter{SD For SHO \& Perturbation Theory}

\section{Schwinger-Dyson For SHO}

We ended last lecture the the derivation of the Schwinger-Dyson equation. We start this lecture by showing how it works for the example of a SHO system. Here the potential is 
\bse 
    V(q) = \frac{m}{2}\omega^2q^2, \qquad \implies \qquad V'(q) = m\omega^2 q,
\ese 
and so we get a nice linear equation for the Schwinger-Dyson formula:
\bse
    \bigg[-\frac{im}{\hbar}\bigg( \omega^2 + \frac{d^2}{dt^2}\bigg) \frac{\del}{\del J(t)} + J(t) \bigg]Z[J] = 0.
\ese 
Using the ansatz $Z[J] = e^{W[J]}$,\footnote{Here we have used the fact that $Z[0]=\braket{0}{0}=1$ for a normalisable theory, otherwise we have to use $Z[J]=Ae^{W[J]}$.} we get 
\bse 
    \frac{im}{\hbar} \bigg( \frac{d^2}{dt^2}+ \omega^2\bigg) \frac{\del W[J]}{\del J(t)} = J(t),
\ese 
so if we take a second functional derivative w.r.t. $J(s)$, we get 
\be 
\label{eqn:WEquation}
    \frac{im}{\hbar} \bigg(\frac{d^2}{dt^2}+\omega^2\bigg) \frac{\del^2 W[J]}{\del J(s)\del J(t)} = \del(s-t).
\ee 
Now we use the definition of the partition function as the sum of $N$-point Green's functions along with the result 
\bse 
    G(t) = \bra{\Omega}q(t) \ket{\Omega} = 0 
\ese
for the SHO we get 
\bse 
    \begin{split}
        W[J] & = \ln Z \\
        & = \ln \bigg( 1 + \frac{1}{2}\int dtds J(t)J(s)G(t,s) + \widetilde{W}[J]\bigg) \\
        & = \frac{1}{2}\int dtds J(t)J(s)G(t,s) + \widetilde{W}[J],
    \end{split}
\ese 
where
\be 
\label{eqn:WTildeEquation}
    \widetilde{W}[J] = \sum_{n=3}^{\infty} \frac{1}{n!}\int dt_1...dt_n J(t_1)...J(t_n) G(t_1,...,t_n) 
\ee 
and we have used
\bse 
    \ln (1+x) = x - \frac{x^2}{2} + \frac{x^3}{3} + ...
\ese 
So we conclude that the $2$-point Green's function is indeed a Green's function with the operator 
\be
\label{eqn:2PointGreensFunction}
    \frac{im}{\hbar} \bigg(\frac{d^2}{dt^2} + \omega^2\bigg) G(t,s) = \del(s-t). 
\ee 
We use our intuition from solving differential equations to suggests the ansatz 
\bse 
    G(t,s) = C e^{-i\omega |s-t|},
\ese 
where we take the modulus as we know that $G(t,s)=G(s,t)$. We can actually go one step further and use the other known fact that 
\be
\label{eqn:2PointExpectationValue}
    G(0,0) = \bra{\Omega}q^2\ket{\Omega} = \frac{\hbar}{2m\omega}
\ee
for the SHO to guess 
\bse 
    G(t,s) = \frac{\hbar}{2m\omega} e^{\pm i\omega |s-t|}.
\ese 

Ok let's check this does indeed satisfy \Cref{eqn:2PointGreensFunction}. We start by rewriting our ansatz as 
\bse 
    G(t,s) = \begin{cases}
        A_+\cos(|s-t|\omega) + B_+\sin(|s-t|\omega) & s>t \\
        A_-\cos(|s-t|\omega) + B_-\sin(|s-t|\omega) & t>s.
    \end{cases}
\ese 
It follows immediately from the symmetry $G(t,s)=G(s,t)$ that we require 
\bse 
    A_+=A_- \qand B_+=B_-,
\ese 
so we can just write 
\bse 
    G(t,s) = A\cos(|s-t|\omega) + B\sin(|s-t|\omega)
\ese 
We get the value of $A$ straight from \Cref{eqn:2PointExpectationValue}, 
\bse 
    A = G(0,0) = \frac{\hbar}{2m\omega}.
\ese 
Now we want the action of the derivative operator in \Cref{eqn:2PointGreensFunction} to give a delta function, which is discontinuous. Therefore, using 
\bse 
    \cos{|x|} =\cos x,
\ese 
which is smooth,\footnote{That is infinitely differentiable, with continuous result.} we know that it is the $\sin$ term that gives us it. Now away from $s=t$ we have a smooth result,  so we only need look at the region $|s-t|\approx 0$, and therefore we can Taylor expand: 
\bse 
    \frac{im}{\hbar} \bigg(\frac{d^2}{dt^2} + \omega^2\bigg)G(t,s) = \frac{im}{\hbar} \bigg(\frac{d^2}{dt^2} + \omega^2\bigg) B|s-t|\omega + \cO(\omega^3)
\ese 
Now we claim that we don't need to consider the $\omega^2$ part of the operator as this will cancel with the $\cO(\omega^3)$ term in the expansion.\footnote{That is we would get a term $\frac{d^2}{dt^2}\big(-\frac{\omega^3}{3!}|s-t|^3\big)$, which cancels the $\omega^2|s-t|\omega$ term.} Then using the fact 
\bse 
    \frac{d}{dx} |x| = \frac{x}{|x|} = \begin{cases} 
        1 & x >0 \\
        -1 & x <0 
        \end{cases} \qquad \implies \qquad \frac{d^2}{dx^2}|x| = 2\del(x),
\ese
we get 
\bse 
    \frac{im}{\hbar} \bigg(\frac{d^2}{dt^2} + \omega^2\bigg)G(t,s) = \frac{im}{\hbar} B\omega 2\del(s-t),
\ese 
so comparing to \Cref{eqn:2PointGreensFunction} we conclude 
\bse 
    B = -\frac{i\hbar}{2m\omega}
\ese 
So altogether we have 
\be
\label{eqn:2PointFunctionSHO}
    G(t,s) = \frac{\hbar}{2m\omega} \big( \cos(|t-s|\omega) - i\sin(|t-s|\omega) \big) = \frac{\hbar}{2m\omega} e^{-i|s-t|\omega}.
\ee 

Now what about $\widetilde{W}[J]$? Well we need it to satisfy 
\bse 
    \frac{im}{\hbar}\bigg( \frac{d^2}{dt^2} + \omega^2\bigg) \frac{\del^2 \widetilde{W}[J]}{\del J(s)\del J(t)} = 0,
\ese 
otherwise we would get another term on the right-hand side of \Cref{eqn:WEquation}. Using the definition \Cref{eqn:WTildeEquation} and the above calculation it's clear we'll get something of the form 
\bse 
    G(t_1,...,t_n) = \big( A_1 \cos(\omega t_1) + B_1\sin(\omega t_1)\big) ... \big( A_n \cos(\omega t_n) + B_n\sin(\omega t_n)\big),
\ese 
but we also require that $G(t_1,...,t_n)$ is totally symmetric, and there is no way to make the above expression totally symmetric apart from in the trivial case 
\bse 
    A_1 = ... = A_n = B_1 = ... = B_n = 0,
\ese
which just leaves us with 
\bse 
    W[J] = \frac{1}{2}\int dt ds J(t_1)J(t_2) G(t_1,t_2), 
\ese 
and so we conclude that the partition function is\footnote{The subscripts $0$ will make sense in a minute.}
\be  
\label{eqn:PartitionFunctionSHO}
    Z_0[J] = \exp\bigg(\frac{1}{2}\int dt_1 dt_2 J(t_1)G_0(t_1,t_2)J(t_2)\bigg),
\ee 
with $G_0(t_1,t_2)$ given by \Cref{eqn:2PointFunctionSHO}.

\section{Perturbation Theory}

In order to make the perturbation calculations we introduce some new notation, as listed below. 
\begin{itemize}
    \item $J_1 = J_{t_1} := J(t)$, 
    \item $(G_0)_{t_1,t_2} := G_0(t_1,t_2)$, 
    \item $J\cdot G_0 \cdot J := \int dt_1dt_2 J_1 (G_0)_{t_1,t_2} J_2$, 
    \item $(J\cdot G_0)_2 := \int dt_1 J_1(G_0)_{t_1,t_2}$, and similarly for $(G_0\cdot J)_1$.
\end{itemize}

Now let's suppose we want to study some system that can be expressed as a perturbation around some known solution, that is the Lagrangian is of the form 
\bse 
    L(q,\dot{q})= L_0(q,\dot{q}) - V(q),
\ese 
where $L_0(q,\dot{q})$ is our known system and $V(q)$ is a `small' perturbation (i.e. the coupling constant that appears in it is small). The partition function for the known system is given by 
\bse 
    Z_0[J] = \cN_0 \int Dq \, \exp\bigg(\frac{i}{\hbar}S_0[q] + \int dt J(t) q(t)\bigg)
\ese 
where $\cN_0$ is some normalisation constant. The full theory then has partition function 
\bse 
    Z[J] = \widetilde{\cN} \int Dq \exp\bigg(\frac{i}{\hbar} S_0[q] + \int dt J(t) q(t) -\frac{i}{\hbar} \int dt V\big(q(t)\big) \bigg).
\ese 
We then employ our clever trick used in the derivation of the Schwinger-Dyson equation and trade the argument $q(t)$ of $V$ for a functional variation w.r.t. $J(t)$ and pull it outside the path integral, giving us 
\bse 
    \begin{split}
        Z[J] & = \exp\bigg[ -\frac{i}{\hbar}\int dt V\bigg(\frac{\del}{\del J(t)}\bigg) \bigg] \widetilde{\cN} \int Dq \exp\bigg(\frac{i}{\hbar} S_0[q] + \int dt J(t)q(t) \bigg) \\
        & = \cN \exp\bigg[ -\frac{i}{\hbar}\int dt V\bigg(\frac{\del}{\del J(t)}\bigg) \bigg] Z_0[J],
    \end{split}
\ese 
where $\cN = \widetilde{\cN}/\cN_0$. Now note that $Z_0[J]$ is independent of the coupling constants (i.e. the parameters in $V(q)$) and so we get a perturbation series by expanding the exponential on the last line above to the required power of coupling.\footnote{If this idea of `power of coupling' doesn't mean anything to you, hopefully it will make sense in a moment, if not reading an introduction to QFT from the second-quantisation approach and Feynman diagrams should clear up any confusion.}

Ok this is all rather abstract, so let's look at an example. 

\bex 
    Let 
    \bse 
        L = \underbrace{\frac{m}{2}\dot{q}^2 - \frac{m}{2}\omega^2q^2}_{L_0 = L_{SHO}} - \underbrace{\frac{\l}{3!}q^3}_{V(q)},
    \ese 
    where we have identified the simply harmonic oscillator as our known system. For clarity with the footnote above, $\l$ here is our coupling constant and we take it to be small. Roughly speaking, it corresponds to the coupling strength for $3$ particles interact (as it comes with $q^3$). In the Taylor expansion the $\l^n$ term will come with $q^{3n}$, and so corresponds to the interaction of $3n$ particles, and this is what we mean by expanding to the relevant power of the coupling constant --- just truncate your expansion to the order of particle interactions you want to consider. 
    
    We have already found the partition function for the SHO, \Cref{eqn:PartitionFunctionSHO}, and so (using the notation introduced at the start of this section) the partition function for the full theory is
    \bse 
        Z[J] = \cN \exp\Bigg[-\frac{i}{\hbar} \frac{\l}{3!} \int dt \bigg(\frac{\del^3}{\del J^3(t)}\bigg)\Bigg] \exp\bigg(\frac{1}{2} J \cdot G_0 \cdot J \bigg).
    \ese 
    Let's just consider coupling to first order, we therefore get 
    \bse 
        Z[J] = \cN \bigg[ 1 - \frac{i}{\hbar} \frac{\l}{3!} \int dt   \bigg(\frac{\del^3}{\del J^3(t)}\bigg) + \cO(\l^2)\bigg]\exp\bigg(\frac{1}{2} J \cdot G_0 \cdot J \bigg).
    \ese 
    We need to find the functional derivatives, and we find this by considering them to be analogous to `regular derivatives' and use the chain rule:\footnote{Note by the symmetry of $G_0$ we have $(J\cdot G_0)_{t_2} = (G_0\cdot J)_{t_1}$, so we can remove the factors of $1/2$ under derivative.}
    \bse 
        \begin{split}
            \frac{\del^3}{\del J^3(t)} \exp\bigg(\frac{1}{2} J \cdot G_0 \cdot J \bigg) & = \frac{\del^2}{\del J^2(t)} \bigg[ (J\cdot G_0)_t \exp\bigg(\frac{1}{2} J \cdot G_0 \cdot J \bigg) \bigg] \\
            & = \frac{\del}{\del J(t)} \bigg[ \Big((G_0)_{t,t} + (J\cdot G_0)^2_t \Big) \exp\bigg(\frac{1}{2} J \cdot G_0 \cdot J \bigg) \bigg] \\
            & = \Big[ (G_0)_{t,t}(J\cdot G_0)_t + 2(G_0)_{t,t}(J\cdot G_0)_t  + (J\cdot G_0)_t^3\Big] \exp\bigg(\frac{1}{2} J \cdot G_0 \cdot J \bigg) \\
            & = \Big[ 3(G_0)_{t,t}(J\cdot G_0)_t + (J\cdot G_0)_t^3\Big] \exp\bigg(\frac{1}{2} J \cdot G_0 \cdot J \bigg).
        \end{split}
    \ese 
    So plugging this into our expression for the partition function, we get 
    \be 
    \label{eqn:PartitionFunctionLambdaCubed}
        Z[J] = \cN \Bigg[ 1 - \frac{i}{\hbar} \l \int dt \bigg( \frac{1}{2}(G_0)_{t,t}(J\cdot G_0)_t + \frac{1}{3!}(J\cdot G_0)_t^3\bigg) \Bigg]\exp\bigg(\frac{1}{2} J \cdot G_0 \cdot J \bigg).
    \ee 
    Finally we can obtain the normalisation coefficient by looking at the $0$-point function (i.e. $Z[0]$) and requiring it be one, This gives 
    \bse 
        1 = \cN \big( 1 + 0\big) e^{0} \qquad \implies \qquad \cN = 1 + \cO(\l^2).
    \ese 
\eex 

\subsection{Feynman Diagrams}

To those familiar with the Hamiltonian approach to QFT, the title of this subsection is probably a huge relief. We can indeed get some Feynman rules from the path integral approach, however one should note that these rules are specific to this case, and so some symbols will appear that you do not see in the second-quantisation Feynman diagrams. We use the above example as a proxy to state the Feynman rules and leave it to the reader to imagine how they adapt to other cases (it is exactly as you would think). 

\mybox{
\begin{center}
	\begin{tabular}{@{} C{4cm} C{4cm} C{4cm} @{}}
		\toprule
		 Part Of Diagram & Maths Expression & Comment \\
		\midrule 
		\btik 
		    \draw[thick] (0,0) -- (2,0); 
		\etik  & \bse G_0(t_1,t_2) \ese  & $t_1/t_2$ are the end points, usually we do not explicitly label them.\footnote{They are always integrated over so just match to the integration variables.} \\
		\btik 
		    \draw[thick] (0,0) circle [radius=0.1cm]; 
		\etik  & \bse \int ds J(s) \ese & The integrand includes the other factors from the diagram. \\
			\btik 
		    \draw[fill=black] (0,0) circle [radius=0.1cm];
		    \draw[thick] (0,0) -- (0,0.75);
		    \draw[thick, rotate around={120:(0,0)}] (0,0) -- (0,0.75);
		    \draw[thick, rotate around={-120:(0,0)}] (0,0) -- (0,0.75);
		\etik  & \bse -\frac{i\l}{\hbar}\int dt  \ese & This is the interaction vertex specifically for our example.\footnote{Note we don't put the $3!$ here, this comes from the symmetry factors, see blow.} This clearly changes depending on the theory.  \\
		\bottomrule
	\end{tabular}
\end{center}
}

As with the second-quantisation approach, we just draw all the possible Feynman diagrams we can at the relevant order of coupling and then obtain the maths expression using the rules above (and also accounting for the symmetry factors). Let's do the example for the previous example. 

\bex 
    We represent the above calculation in Feynman diagrams as
    \begin{center}
        \btik 
            \node at (0,0) {\Large{$\frac{1}{\cN}Z[J] = \Bigg[ 1 \, +$}};
            \draw[thick] (2.25,0) circle [radius=0.5cm];
            \draw[fill=black] (2.75,0) circle [radius=0.1cm]; 
            \draw[thick] (2.75,0) -- (3.75,0);
            \draw[thick, fill=white] (3.75,0) circle [radius=0.1cm];
            \node at (4.4,0) {\Large{$+$}};
            \draw[fill=black] (5.5,0) circle [radius=0.1cm];
		    \draw[thick] (5.5,0) -- (5.5,0.75);
		    \draw[thick, fill=white] (5.5,0.75) circle [radius=0.1cm];
		    \draw[thick, rotate around={120:(5.5,0)}] (5.5,0) -- (5.5,0.75);
		    \draw[thick, fill=white, rotate around={120:(5.5,0)}] (5.5,0.75) circle [radius=0.1cm];
		    \draw[thick, rotate around={-120:(5.5,0)}] (5.5,0) -- (5.5,0.75);
		    \draw[thick, fill=white, rotate around={-120:(5.5,0)}] (5.5,0.75) circle [radius=0.1cm];
		    \node at (8,0) {\Large{$+ \, \cO(\l^2) \Bigg] \exp\bigg($}};
		    \draw[thick] (9.75,0) -- (10.75,0);
		    \draw[thick, fill=white] (9.75,0) circle [radius=0.1cm];
		    \draw[thick, fill=white] (10.75,0) circle [radius=0.1cm];
		    \node at (11,0) {\Large{$\bigg)$}};
        \etik  
    \end{center}
    We can then work out the symmetry factors by looking at the diagrams:
    \begin{center}
	\begin{tabular}{@{} C{4cm} C{4cm} C{4cm} @{}}
		\toprule
		 Diagram & Symmetry Factor & Why \\
		\midrule 
		\btik 
		    \draw[thick] (2.25,0) circle [radius=0.5cm];
            \draw[fill=black] (2.75,0) circle [radius=0.1cm]; 
            \draw[thick] (2.75,0) -- (3.75,0);
            \draw[thick, fill=white] (3.75,0) circle [radius=0.1cm];
		\etik & \bse \frac{1}{2} \ese & We can reflect the loop about the horizontal axis. \\
		\btik 
		    \draw[fill=black] (5.5,0) circle [radius=0.1cm];
		    \draw[thick] (5.5,0) -- (5.5,0.75);
		    \draw[thick, fill=white] (5.5,0.75) circle [radius=0.1cm];
		    \draw[thick, rotate around={120:(5.5,0)}] (5.5,0) -- (5.5,0.75);
		    \draw[thick, fill=white, rotate around={120:(5.5,0)}] (5.5,0.75) circle [radius=0.1cm];
		    \draw[thick, rotate around={-120:(5.5,0)}] (5.5,0) -- (5.5,0.75);
		    \draw[thick, fill=white, rotate around={-120:(5.5,0)}] (5.5,0.75) circle [radius=0.1cm];
		\etik & \bse \frac{1}{3!} \ese & Permutations of the three lines. \\
		\btik 
		    \draw[thick] (9.75,0) -- (10.75,0);
		    \draw[thick, fill=white] (9.75,0) circle [radius=0.1cm];
		    \draw[thick, fill=white] (10.75,0) circle [radius=0.1cm];
		\etik & \bse \frac{1}{2} \ese & Can reflect in a vertical line down the middle. \\
		\bottomrule
	\end{tabular}
\end{center}
\eex 

\bbox 
    Convince yourself that the diagram given above does indeed give us \Cref{eqn:PartitionFunctionLambdaCubed}. 
\ebox 

\subsection{Green's Functions From Feynman Diagrams}

Recall that we can obtain the $N$-point Green's function by taking functional derivatives w.r.t. $J(t)$ of the partition function, \Cref{eqn:GreensFunctionsFromParition}:
\bse 
    G(t_1,..,t_N) = \frac{\del^N Z[J]}{\del J(t_1) ... \del J(t_N)} \bigg|_{J=0}.
\ese
We now want to see how to do this in terms of our diagrams. Well the hollow circles represent integrated factors of $J(s)$ in the diagram, so our functional varitation will remove this circle and replace it with the derivative $J$s argument, i.e. we use, for example,
\bse 
    \frac{\del^2}{\del J_1 J_2} \bigg( \frac{1}{2} (J\cdot G_0\cdot J)_{s,r} \bigg) =  (G_0)_{t_1,t_2}
\ese 
to replace
\begin{center}
    \btik 
        \node at (-1,0) {\Large{$\frac{\del^2}{\del J_1 \del J_2}\Big($}};
        \draw[thick] (0,0) -- (2,0);
        \draw[thick, fill=white] (0,0) circle [radius=0.1cm];
        \draw[thick, fill=white] (2,0) circle [radius=0.1cm]; 
        \node at (2.75,0) {\Large{$\Big) \, = $}};
        \draw[thick] (4,0) -- (6,0);
        \node at (4,-0.3) {$t_1$};
        \node at (6,-0.3) {$t_2$};
    \etik 
\end{center}

Note that it is only the terms that have exactly $N$ hollow circles in it that survive this process: if there's less then there's a derivative that has nothing to act on; if there's more then setting $J=0$ corresponds to saying diagrams with a hollow circle left vanish. 

\br 
    Note that our differentiated diagrams need to respect the symmetry property of the maths. For example, we know that $(G_0)_{t_1,t_2} = (G_0)_{t_2,t_1}$ and so we must have 
    \begin{center}
        \btik 
            \draw[thick] (0,0) -- (2,0);
            \node at (0,-0.3) {$t_1$};
            \node at (2,-0.3) {$t_2$};
            \node at (3,0) {\Large{$=$}};
            \draw[thick] (4,0) -- (6,0);
            \node at (4,-0.3) {$t_2$};
            \node at (6,-0.3) {$t_1$};
        \etik 
    \end{center}
    Note we don't get factors from this symmetry though, because the diagrams come with that same factor divided. For example  
    \begin{center}
        \btik 
            \node at (-1.4,0) {\Large{$\frac{\del}{\del J_1}\Bigg($}};
            \draw[fill=black] (0,0) circle [radius=0.1cm];
		    \draw[thick] (0,0) -- (0,0.75);
		    \draw[thick, fill=white] (0,0.75) circle [radius=0.1cm];
		    \draw[thick, rotate around={120:(0,0)}] (0,0) -- (0,0.75);
		    \draw[thick, fill=white, rotate around={120:(0,0)}] (0,0.75) circle [radius=0.1cm];
		    \draw[thick, rotate around={-120:(0,0)}] (0,0) -- (0,0.75);
		    \draw[thick, fill=white, rotate around={-120:(0,0)}] (0,0.75) circle [radius=0.1cm];
		    \node at (1.4,0) {\Large{$\Bigg) \, =$}}; 
		    \begin{scope}[xshift=3cm]
		        \draw[fill=black] (0,0) circle [radius=0.1cm];
    		    \draw[thick] (0,0) -- (0,0.75);
    		    \node at (0,1) {$t_1$};
    		    \draw[thick, rotate around={120:(0,0)}] (0,0) -- (0,0.75);
    		    \draw[thick, fill=white, rotate around={120:(0,0)}] (0,0.75) circle [radius=0.1cm];
    		    \draw[thick, rotate around={-120:(0,0)}] (0,0) -- (0,0.75);
    		    \draw[thick, fill=white, rotate around={-120:(0,0)}] (0,0.75) circle [radius=0.1cm];
		    \end{scope}
        \etik 
    \end{center}
    \textit{without} a factor of $3!$ at the front. That is, there are $3!$ different ways to get the right-hand side, but the symmetry factor of this diagram is exactly $1/3!$.
\er


%\br 
%    We need to take the symmetry factors into account when taking derivatives of the diagrams. By which we mean, imagine temporarily colouring the hollow circles to track them, then write out all the permutations and then take the derivative always removing the circle in the same position. If we then remove the colouring, all the diagrams add up and we might try stick a prefactor in front of the derivative. However we have to remember that we also divide by the symmetry factor in the diagrams, and this will exactly cancel this prefactor. Perhaps an example will help see this: imagine differentiating this \textcolor{red}{Ask Douglas}
%    \begin{center}
%        \btik 
%            \draw[thick] (0,0) -- (2,0);
%            \draw[thick, fill=white] (0,0) circle [radius=0.1cm];
%            \draw[thick, fill=white] (2,0) circle [radius=0.1cm];
%            \node at (3,0) {\Large{$= \, \frac{1}{2!}\bigg($}};
%            \draw[thick] (4,0) -- (6,0);
%            \draw[thick, blue, fill=white] (4,0) circle [radius=0.1cm];
%            \draw[thick, red, fill=white] (6,0) circle [radius=0.1cm];
%            \node at (6.5,0) {\Large{$+$}};
%            \draw[thick] (7,0) -- (9,0);
%            \draw[thick, red, fill=white] (7,0) circle [radius=0.1cm];
%            \draw[thick, blue, fill=white] (9,0) circle [radius=0.1cm];
%            \node at (9.5,0) {\Large{$\bigg)$}};
%        \etik 
%    \end{center}
%\er 

\bex 
    We can find the $2$-point function for out $q^3$ theory. To first order in $\l$, all the diagrams outside the exponential vanish (as none have exactly $2$ hollow circles), and so we are just left with the terms that are pulled down from the action on the exponential. It is clear\footnote{And hence I'm saving myself the pain of Tikzing it all.} that what we are left with at the end is simply 
    \begin{center}
        \btik 
            \draw[thick] (0,0) -- (2,0);
            \node at (0,-0.3) {$t_1$};
            \node at (2,-0.3) {$t_2$};
        \etik  
    \end{center}
\eex 

We now note that from the diagrams we can obtain terms that do not appear in the expansion of \Cref{eqn:PartitionFunctionLambdaCubed}. For example, the following diagram is a valid diagram at first order in $\l$
\begin{center}
    \btik 
        \draw[thick] (0,0) -- (1,0);
        \draw[thick] (-0.5,0) circle [radius=0.5cm];
        \draw[fill=black] (0,0) circle [radius=0.1cm];
        \draw[thick, fill=white] (1,0) circle [radius=0.1cm];
        \draw[thick] (1.3,0) -- (2.3,0);
        \draw[thick, fill=white] (1.3,0) circle [radius=0.1cm];
        \draw[thick, fill=white] (2.3,0) circle [radius=0.1cm];
    \etik  
\end{center}
\noindent where we note this is \textit{not} the sum of two diagrams, but is collectively one diagram. We see that this is essentially two of the previous diagrams `put next to each other' in a disconnected way. This diagram will give a non-zero contribution to the $3$-point Green's function, and it is easy to convince yourself (exercise below) that the full $3$-point function is given by 
\begin{center}
    \btik 
        \begin{scope}[xshift=-4cm]
            \draw[thick] (0,0) -- (1,0);
            \draw[thick] (-0.5,0) circle [radius=0.5cm];
            \draw[fill=black] (0,0) circle [radius=0.1cm];
            \node at (1,-0.3) {$t_1$};
            \draw[thick] (1.3,0) -- (2.3,0);
            \node at (1.4,-0.3) {$t_2$};
            \node at (2.3,-0.3) {$t_3$};
            \node at (2.7,0) {\Large{$+$}};
        \end{scope}
        \begin{scope}[]
            \draw[thick] (0,0) -- (1,0);
            \draw[thick] (-0.5,0) circle [radius=0.5cm];
            \draw[fill=black] (0,0) circle [radius=0.1cm];
            \node at (1,-0.3) {$t_2$};
            \draw[thick] (1.3,0) -- (2.3,0);
            \node at (1.4,-0.3) {$t_1$};
            \node at (2.3,-0.3) {$t_3$};
            \node at (2.7,0) {\Large{$+$}};
        \end{scope}
        \begin{scope}[xshift=4cm]
            \draw[thick] (0,0) -- (1,0);
            \draw[thick] (-0.5,0) circle [radius=0.5cm];
            \draw[fill=black] (0,0) circle [radius=0.1cm];
            \node at (1,-0.3) {$t_3$};
            \draw[thick] (1.3,0) -- (2.3,0);
            \node at (1.4,-0.3) {$t_2$};
            \node at (2.3,-0.3) {$t_1$};
            \node at (2.7,0) {\Large{$+$}};
        \end{scope}
        \begin{scope}[xshift=8cm]
            \draw[fill=black] (0,0) circle [radius=0.1cm];
		    \draw[thick] (0,0) -- (0,0.75);
		    \node at (0,1) {$t_1$};
		    \draw[thick, rotate around={120:(0,0)}] (0,0) -- (0,0.75);
		    \node at (1,-0.5) {$t_2$};
		    \draw[thick, rotate around={-120:(0,0)}] (0,0) -- (0,0.75);
		    \node at (-0.9,-0.5) {$t_3$};
        \end{scope}
    \etik 
\end{center}

\bbox 
    Convince yourself that the above is indeed the $3$-point function for the $q^3$ theory. 
\ebox 

\subsection{Vacuum Bubbles}

Now note that if we go to order $\l^2$ then we will get diagrams that look like 
\begin{center}
    \btik
        \begin{scope}[xshift=-3cm]
            \draw[thick] (0,0) -- (1,0);
            \draw[thick] (-0.5,0) circle [radius=0.5cm];
            \draw[fill=black] (0,0) circle [radius=0.1cm];
            \draw[fill=black] (1,0) circle [radius=0.1cm];
            \draw[thick] (1.5,0) circle [radius=0.5cm];
        \end{scope}
        \node at (0.2,0) {and};
        \begin{scope}[xshift=3cm]
            \draw[thick] (-1,0) -- (1,0);
            \draw[thick] (0,0) circle [radius=1cm];
            \draw[fill=black] (-1,0) circle [radius=0.1cm];
            \draw[fill=black] (1,0) circle [radius=0.1cm];
        \end{scope}
    \etik 
\end{center}
These diagrams have no external legs (i.e. there are no hollow circles), and we refer to diagrams of this kind as \textit{vacuum bubbles}. Our partition function does not contain vacuum bubbles and so we need to  `divide the diagrams out'. What do we mean by this? Well we note that for every one of these vacuum bubble diagrams we also get the same diagram but now with a $(J\cdot G_0 \cdot J)$ term too, e.g.  
\begin{center}
    \btik 
        \draw[thick] (0,0) -- (1,0);
        \draw[thick] (-0.5,0) circle [radius=0.5cm];
        \draw[fill=black] (0,0) circle [radius=0.1cm];
        \draw[fill=black] (1,0) circle [radius=0.1cm];
        \draw[thick] (1.5,0) circle [radius=0.5cm];
        \draw[thick] (2.3,0) -- (3.3,0);
        \draw[thick,fill=white] (2.3,0) circle [radius=0.1cm];
        \draw[thick,fill=white] (3.3,0) circle [radius=0.1cm];
    \etik
\end{center}
We therefore `factor out' the vacuum bubble parts and divide by the purely vacuum bubble diagrams to cancel them. For example, the $2$-point function for the $q^3$ theory at order $\l^3$ is given by the division of diagrams
\begin{center}
    \btik 
        \begin{scope}
            \draw[thick] (0,0) -- (1,0);
            \node at (0,-0.3) {$t_1$};
            \node at (1,-0.3) {$t_2$};
        \end{scope}
        \begin{scope}[xshift=3.75cm]
            \node at (-1.75,0) {\Large{$\Bigg[ 1 \, +$}};
            %
            \draw[thick] (0,0) -- (1,0);
            \draw[thick] (-0.5,0) circle [radius=0.5cm];
            \draw[fill=black] (0,0) circle [radius=0.1cm];
            \draw[fill=black] (1,0) circle [radius=0.1cm];
            \draw[thick] (1.5,0) circle [radius=0.5cm];
            %
            \node at (2.3,0) {\Large{$+$}};
            %
            \draw[thick] (2.7,0) -- (3.7,0);
            \draw[thick] (3.2,0) circle [radius=0.5cm];
            \draw[fill=black] (2.7,0) circle [radius=0.1cm];
            \draw[fill=black] (3.7,0) circle [radius=0.1cm];
            %
            \node at (4.2,0) {\Large{$\Bigg] \, +$}};
        \end{scope}
        \begin{scope}[xshift=9cm]
            \draw[thick] (0,-0.7) -- (0,0);
            \draw[thick] (0,0.4) circle [radius=0.4cm];
            \draw[fill=black] (0,0) circle [radius=0.1cm];
            \node at (0,-1) {$t_1$};
            \draw[thick] (1,-0.7) -- (1,0);
            \draw[thick] (1,0.4) circle [radius=0.4cm];
            \draw[fill=black] (1,0) circle [radius=0.1cm];
            \node at (1,-1) {$t_2$};
            \node at (1.6,0) {\Large{$+$}};
        \end{scope}
        \begin{scope}[xshift=12cm]
            \draw[thick] (-1,0) -- (-0.5,0);
            \node at (-0.9,-0.3) {$t_1$};
            \draw[fill=black] (-0.5,0) circle [radius=0.1cm];
            \draw[thick] (0,0) circle [radius=0.5cm];
            \draw[fill=black] (0.5,0) circle [radius=0.1cm];
            \draw[thick] (0.5,0) -- (1,0);
            \node at (1,-0.3) {$t_2$};
            \node at (1.4,0) {\Large{$+$}};
        \end{scope}
        \begin{scope}[xshift=14cm]
            \draw[thick] (0,-0.5) -- (1.2,-0.5);
            \node at (0.1,-0.8) {$t_1$};
            \node at (1.2,-0.8) {$t_2$};
            \draw[fill=black] (0.6,-0.5) circle [radius=0.1cm];
            \draw[thick] (0.6,-0.5) -- (0.6,0);
            \draw[thick] (0.6,0.4) circle [radius=0.4cm];
            \draw[fill=black] (0.6,0) circle [radius=0.1cm];
        \end{scope}
        \draw[ultra thick] (0,-1.5) -- (15.5,-1.5);
        \begin{scope}[xshift=6.5cm,yshift=-2.5cm]
            \node at (-1.75,0) {\Large{$\Bigg[ 1 \, +$}};
            %
            \draw[thick] (0,0) -- (1,0);
            \draw[thick] (-0.5,0) circle [radius=0.5cm];
            \draw[fill=black] (0,0) circle [radius=0.1cm];
            \draw[fill=black] (1,0) circle [radius=0.1cm];
            \draw[thick] (1.5,0) circle [radius=0.5cm];
            %
            \node at (2.3,0) {\Large{$+$}};
            %
            \draw[thick] (2.7,0) -- (3.7,0);
            \draw[thick] (3.2,0) circle [radius=0.5cm];
            \draw[fill=black] (2.7,0) circle [radius=0.1cm];
            \draw[fill=black] (3.7,0) circle [radius=0.1cm];
            %
            \node at (4.2,0) {\Large{$\Bigg]$}};
        \end{scope}
    \etik  
\end{center}
\noindent which after the division just leaves 
\begin{center}
    \btik 
        \begin{scope}
            \draw[thick] (0,0) -- (1,0);
            \node at (0,-0.3) {$t_1$};
            \node at (1,-0.3) {$t_2$};
            \node at (1.5,0) {\Large{$+$}};
        \end{scope}
        \begin{scope}[xshift=2.5cm]
            \draw[thick] (0,-0.7) -- (0,0);
            \draw[thick] (0,0.4) circle [radius=0.4cm];
            \draw[fill=black] (0,0) circle [radius=0.1cm];
            \node at (0,-1) {$t_1$};
            \draw[thick] (1,-0.7) -- (1,0);
            \draw[thick] (1,0.4) circle [radius=0.4cm];
            \draw[fill=black] (1,0) circle [radius=0.1cm];
            \node at (1,-1) {$t_2$};
            \node at (1.8,0) {\Large{$+$}};
        \end{scope}
        \begin{scope}[xshift=6cm]
            \draw[thick] (-1,0) -- (-0.5,0);
            \node at (-0.9,-0.3) {$t_1$};
            \draw[fill=black] (-0.5,0) circle [radius=0.1cm];
            \draw[thick] (0,0) circle [radius=0.5cm];
            \draw[fill=black] (0.5,0) circle [radius=0.1cm];
            \draw[thick] (0.5,0) -- (1,0);
            \node at (1,-0.3) {$t_2$};
            \node at (1.5,0) {\Large{$+$}};
        \end{scope}
        \begin{scope}[xshift=8cm]
            \draw[thick] (0,-0.5) -- (1.2,-0.5);
            \node at (0.1,-0.8) {$t_1$};
            \node at (1.2,-0.8) {$t_2$};
            \draw[fill=black] (0.6,-0.5) circle [radius=0.1cm];
            \draw[thick] (0.6,-0.5) -- (0.6,0);
            \draw[thick] (0.6,0.4) circle [radius=0.4cm];
            \draw[fill=black] (0.6,0) circle [radius=0.1cm];
        \end{scope}
    \etik 
\end{center}

This gives us the final general conclusion 
\mybox{
\begin{center}
    $N$-point function is given by sum of all diagrams with exactly $N$ external legs but no vacuum bubble contributions. 
\end{center}
}

\bbox 
    Convince yourself that there are no other contributions to the $2$-point function at order $\l^2$. 
\ebox 

\section{Semi-Classical Expansion}

We conclude the course with a brief discussion of the semi-classical approximation to the path integral. We obtain this by Taylor expanding the action about a classical path $q_{cl}$:
\bse 
    S[q_{cl}+\eta] = S[q_{cl}] + \int dt_1 \eta(t_1) \frac{\del S}{\del q(t_1)} [q_{cl}] + \frac{1}{2!}\int dt_1 dt_2 \eta(t_1)\eta(t_2) \frac{\del^2 S}{\del q(t_1) \del q(t_2)} [q_{cl}] + ...,
\ese 
now we note that the term with a single integral vanishes as the variation of the action gives the Euler-Lagrange equations and the classical path minimises these. Now let's consider rescalling 
\bse 
    \eta \to \sqrt{\hbar}\eta,
\ese 
and noting, as before,
\bse 
    D(q_{cl}+\eta) = D\eta,
\ese
then our Feynman kernel becomes 
\bse 
    \begin{split}
        K(q_F,T;q_I,0) & = \int Dq \, \exp\bigg(\frac{i}{\hbar} S[q_{cl}+\eta]\bigg) \\
        & \propto e^{\frac{i}{\hbar}S[q_{cl}]}\int D\eta \, \exp\Bigg[ \frac{i}{2!} \int dt_1 dt_2 \eta(t_1)\eta(t_2) \frac{\del^2 S[q_{cl}]}{\del q(t_1) \del q(t_2)} + ... \\
        & \qquad + \frac{i}{n!} \hbar^{\frac{n}{2}-1} \int dt_1 ... dt_n \eta(t_1)...\eta(t_n) \frac{\del^n S[q_{cl}]}{\del q(t_1)... \del q(t_n)}  + ... \Bigg] \\
        & = e^{\frac{i}{\hbar}S[q_{cl}]}\int D\eta \, \exp\Bigg[ \frac{i}{2!} \int dt_1 dt_2 \eta(t_1)\eta(t_2) \frac{\del^2 S[q_{cl}]}{\del q(t_1) \del q(t_2)}\Bigg]\Big( 1 +\cO\big(\sqrt{\hbar}\big) \Big).
    \end{split}
\ese 
Now obviously the prefactor is the classical contribution, and we see that the leading order quantum contribution is a Gaussian path integral in $\eta$, i.e. its of the form
\bse 
    \int D \eta \, \exp\big( i\eta\cdot \Omega \cdot \eta\big) \propto \big(\det (-i\Omega)\big)^{-1/2},
\ese 
where\footnote{Extra exercise: prove this formula holds.}
\bse 
    \begin{split}
        \Omega_{t_1,t_2} & = \frac{1}{2} \frac{\del^2 S[q_{cl}]}{\del q(t_1) \del q(t_2)} \\
        & = -\frac{1}{2}\bigg( m\frac{d^2}{dt_1^2} + V''\big(q_{cl}\big)\bigg) \del(t_1-t_2)
    \end{split}
\ese 
where we have assumed the Lagrangian has the usual kinetic term (i.e. $\frac{m}{2}\dot{q}^2$). $\det(-i\Omega)$ is the product of the eigenvalues of the operator
\bse 
    \frac{i}{2}\bigg( m\frac{d^2}{dt_1^2} + V''\big(q_{cl}\big)\bigg)
\ese 
acting on $\eta(t)$ subject to the constraints $\eta(t_I)=\eta(t_F)=0$. The semiclassical approximation is when we ignore the contributions from the $\cO(\sqrt{\hbar})$ terms.

\br 
    If the action is quadratic in $q$ (e.g. the SHO) then the semiclassical solution is exact.
\er 

\br 
    If the classical path is not unique we must take a sum over all the different classical solutions.
\er 

\bex 
    Let's look at the SHO as an example. As just remarked, this will give us an exact result. Here we have 
    \bse 
        V''(q_{cl}) = m\omega^2, \qquad \implies \qquad \Omega_{t_1,t_2} = -\frac{m}{2} \bigg(\frac{d^2}{dt_1^2} + \omega^2\bigg)\del(t_1-t_2).
    \ese 
    The eigenvectors of the operator here are 
    \bse 
        \eta(t) = A \cos(\widetilde{\omega}t) + B \sin(\widetilde{\omega}t),
    \ese 
    with eigenvalues 
    \bse 
        \l = \omega^2 -\widetilde{\omega}^2.
    \ese 
    If we now impose the condition $\eta(t_I)=0=\eta(t_F)$ with $t_I=0$, $t_F=T$, we conclude 
    \bse 
        \widetilde{\omega} = \frac{n\pi}{T}, \qquad n \in \N,
    \ese 
    and so 
    \bse 
        \l = \omega^2 - \bigg(\frac{n\pi}{T}\bigg)^2 \propto 1 - \bigg(\frac{\omega T}{n\pi}\bigg)^2.
    \ese
    We can therefore conclude 
    \bse 
        \big(\det(-i\Omega)\big)^{-1/2} \propto \prod_{n=1}^{\infty} \bigg[ 1 - \bigg(\frac{\omega T}{n\pi}\bigg)^2\bigg]^{-1/2} = \sqrt{\frac{\omega T}{\sin(\omega T)}}.
    \ese 
\eex 

% -------------------------------------------------------------------
% Bibliography/Further Readings
% -------------------------------------------------------------------

\chapter*{Useful Texts \& Further Readings}

\section*{QFT}
\begin{itemize}
    \item Zee, Anthony. \textit{Quantum Field Theory In A Nutshell}. Vol. 7. Princeton University Press, 2010.
\end{itemize}




%\bibliographystyle{agsm} 
%\bibliography{mybibliography} 
%\printbibliography[heading=bibintoc]


% -------------------------------------------------------------------
% Appendices
% -------------------------------------------------------------------

%\begin{appendices}
%\chapter{An Appendix of Some Kind}

\lipsum[1-3]  % Replace with your text

%\end{appendices}

\end{document}
